% =============================================================================
% L05 (Classes 5+6) --- textbook exercises with solutions. Built from the MPE list template
% (~/.claude/commands/templates/lista-exam.tex).
%
% TWO MODES (one source file):
%   \solucoesfalse  -> blank LIST (space to solve)
%   \solucoestrue   -> ANSWER KEY (same file)
%
% Compile (copyable PDF, global lmodern+cmap rule):
%   pdflatex lista-L05-book-exercises.tex  (run 2x) ; check: pdffonts ... (all Type 1)
% Numbers: python Listas/L05-book-exercises-check.py reproduces and asserts every number.
% =============================================================================
\documentclass[11pt,a4paper]{article}
\usepackage[utf8]{inputenc}
\usepackage[T1]{fontenc}
\usepackage{lmodern}   % Type 1 vector fonts (replaces PK bitmaps)
\usepackage{cmap}      % ToUnicode CMap -> copyable text and accents
\usepackage[english]{babel}
\usepackage{amsmath,amssymb,amsfonts}
\usepackage[a4paper,margin=2.2cm]{geometry}
\usepackage{enumitem}
\usepackage{fancyhdr}
\usepackage{comment}
% One cut file per job: parallel builds in one folder would otherwise share comment.cut and
% swap solutions. Redefine DefaultCutFileName; overriding CommentCutFile alone is reset by comment.sty.
\def\DefaultCutFileName{\def\CommentCutFile{\jobname.cut}}
\usepackage{needspace}
\usepackage{xcolor}
\usepackage{microtype}
\usepackage{booktabs}
\usepackage{graphicx}

% ====================== CONFIGURATION ======================
\newif\ifsolucoes
\solucoesfalse                 % <<<< \solucoestrue for the ANSWER KEY; \solucoesfalse for the LIST

\newcommand{\disciplina}{Econometria Avançada --- MPE (Insper)}
\newcommand{\listatema}{L05 --- Introduction to causal inference (Classes 5+6): textbook exercises}
\newcommand{\espacoresol}{4.5cm}   % height of the solving space in the blank list

% ====================== MACHINERY ======================
\ifsolucoes\includecomment{sol}\else\excludecomment{sol}\fi

\pagestyle{fancy}\fancyhf{}
\lhead{\footnotesize\disciplina}
\rhead{\footnotesize List L05\ifsolucoes{} --- Solutions\fi}
\cfoot{\footnotesize\thepage}
\renewcommand{\headrulewidth}{0.3pt}
\setlength{\headheight}{14pt}

\setlist{topsep=2pt,itemsep=2pt}
\newcounter{questao}
% Equations are numbered per question, (Q<question>.<n>), so they never collide with the
% books' own numbers, which keep their printed tags, e.g. (11.4) or (13.9).
\renewcommand{\theequation}{Q\thequestao.\arabic{equation}}

% Block header
\newcommand{\bloco}[3]{\bigskip\noindent
  {\large\bfseries\color{blue!55!black}Block #1 --- #2}\hfill\textit{(#3)}\par
  \noindent\rule{\linewidth}{0.3pt}\smallskip}

% Question (automatic numbering; resets the equation counter)
\newcommand{\questao}[1]{\refstepcounter{questao}\setcounter{equation}{0}\medskip\needspace{5\baselineskip}\noindent
  \textbf{Question \thequestao.}~#1\par}
% Source line that opens every question (estilo-avaliacao §5, book questions)
\newcommand{\fonte}[1]{\par\noindent{\small\textit{Source:} #1}\par\smallskip}

% Notation list that opens every solution (estilo-avaliacao §4, "Equations, notation and layout").
\newenvironment{notacao}{\par\smallskip\noindent\textbf{Notation.}\begin{itemize}[leftmargin=1.4em,topsep=1pt,itemsep=0pt]\small}{\end{itemize}}
\allowdisplaybreaks

% Solving space in the blank list
\newcommand{\resolespaco}{\ifsolucoes\else\par\textit{Solution:}\vspace{\espacoresol}\par\fi}

% Solution step (use inside sol)
\newcommand{\passo}[2]{\par\smallskip\noindent\textbf{Step --- #1.}~#2}
\newcommand{\intuicao}[1]{\par\smallskip\noindent\textit{Intuition.}~#1}
\newcommand{\refbib}[1]{\par\smallskip\noindent{\footnotesize\color{black!60}\textbf{Ref:} #1}}

% =============================================================================
\begin{document}
\begin{center}
  {\Large\bfseries \disciplina}\\[2pt]
  {\large Exercise list: \listatema}\ifsolucoes\\[2pt]\textit{Solutions}\fi
\end{center}
\vspace{2pt}

{\footnotesize\noindent\textbf{Instructions.} These are the \emph{Practice} items of Classes 5+6 in
\texttt{Map/class-plan.md}. Every question opens with a source line. Book statements (Wooldridge,
Greenberg) are reproduced verbatim from the book PDF. \emph{Mostly Harmless Econometrics} has no
end-of-chapter exercises: for each of its worked examples the table is reproduced verbatim and the
question is original. Computer exercises are solved in Python (\texttt{statsmodels},
\texttt{linearmodels}); \texttt{Listas/L05-book-exercises-check.py} reproduces and asserts every number
in the solutions. Show your steps.\par}

% =============================================================================
\bloco{1}{MHE worked examples}{MHE Ch.~2--6 \,\textbullet\, Questions 1--6}

% ---------------------------------------------------------------- Q1
\questao{The STAR experiment: balance, selection bias and covariates (MHE Tables 2.2.1--2.2.2).}
\fonte{Angrist and Pischke (2008), \textit{Mostly Harmless Econometrics: An Empiricist's Companion},
\S2.2, Tables 2.2.1 and 2.2.2, printed pp.~14--15. The source has no exercise; the statement is
original. The table numbers are reproduced verbatim from the PDF (table notes shortened).}

\begin{center}\small
\textit{Table 2.2.1: Comparison of treatment and control characteristics in the Tennessee STAR experiment}\\[2pt]
\begin{tabular}{llrrrr}
\toprule
\multicolumn{6}{c}{Students who entered STAR in kindergarten}\\
& Variable & Small & Regular & Regular/Aide & Joint $P$-value\\
\midrule
1. & Free lunch & .47 & .48 & .50 & .09\\
2. & White/Asian & .68 & .67 & .66 & .26\\
3. & Age in 1985 & 5.44 & 5.43 & 5.42 & .32\\
4. & Attrition rate & .49 & .52 & .53 & .02\\
5. & Class size in kindergarten & 15.10 & 22.40 & 22.80 & .00\\
6. & Percentile score in kindergarten & 54.70 & 48.90 & 50.00 & .00\\
\bottomrule
\end{tabular}\\[2pt]
\parbox{0.9\linewidth}{\footnotesize Notes (excerpt): ``Adapted from Krueger (1999), Table 1. The table shows
means of variables by treatment status. The $P$-value in the last column is for the $F$-test of
equality of variable means across all three groups.''}

\medskip
\textit{Table 2.2.2: Experimental estimates of the effect of class-size assignment on test scores}\\[2pt]
\begin{tabular}{lcccc}
\toprule
Explanatory variable & (1) & (2) & (3) & (4)\\
\midrule
Small class & 4.82 & 5.37 & 5.36 & 5.37\\
 & (2.19) & (1.26) & (1.21) & (1.19)\\
Regular/aide class & .12 & .29 & .53 & .31\\
 & (2.23) & (1.13) & (1.09) & (1.07)\\
White/Asian (1 = yes) & -- & -- & 8.35 & 8.44\\
 & & & (1.35) & (1.36)\\
Girl (1 = yes) & -- & -- & 4.48 & 4.39\\
 & & & (.63) & (.63)\\
Free lunch (1 = yes) & -- & -- & -13.15 & -13.07\\
 & & & (.77) & (.77)\\
White teacher & -- & -- & -- & -.57\\
 & & & & (2.10)\\
Teacher experience & -- & -- & -- & .26\\
 & & & & (.10)\\
Master's degree & -- & -- & -- & -0.51\\
 & & & & (1.06)\\
School fixed effects & No & Yes & Yes & Yes\\
R$^2$ & .01 & .25 & .31 & .31\\
\bottomrule
\end{tabular}\\[2pt]
\parbox{0.9\linewidth}{\footnotesize Note (excerpt): ``The dependent variable is the Stanford Achievement Test
percentile score. Robust standard errors that allow for correlated residuals within classes are shown
in parentheses. The sample size is 5681.''}
\end{center}

\begin{enumerate}[label=(\alph*)]
\item What null hypothesis does the ``Joint $P$-value'' column test? Using rows 1--3 of Table 2.2.1,
did randomization balance the groups at the 5\% level? Why should rows 4--6 not be read as balance checks?
\item Write the observed difference in mean scores between small and regular classes as the effect on
the treated plus selection bias, and explain why random assignment makes the coefficient on
\emph{Small class} an estimate of the average causal effect.
\item Using column (1) of Table 2.2.2, compute the $t$ statistic and a 95\% confidence interval for the
small-class effect; repeat with column (4). Is the regular/aide effect in column (1) significant?
\item From column (2) to column (3) the small-class coefficient barely moves while its standard error
falls from 1.26 to 1.21. Explain both facts, and show that the fall in the standard error is about what
the rise in $R^2$ from .25 to .31 predicts.
\item The text says the effect is about $.2\sigma$, where $\sigma$ is the standard deviation of the
kindergarten percentile score. What $\sigma$ does column (4) imply? Is it plausible for a percentile scale?
\end{enumerate}
\resolespaco
\begin{sol}
\begin{notacao}
\item $Y_i$: percentile score of pupil $i$ (0--100); $D_i\in\{0,1\}$: assigned to a small class;
$Y_{1i},Y_{0i}$: potential scores with and without a small class.
\item $\hat\rho$: estimated coefficient on \emph{Small class}; $\widehat{se}$: its standard error;
$t=\hat\rho/\widehat{se}$; $X_i$: covariates (race, sex, free lunch, teacher traits).
\item $R^2_k$: $R^2$ of column $k$; $\sigma$: standard deviation of the percentile score ($\sigma>0$).
\end{notacao}

\passo{(a) balance}{For each variable the test is $H_0:\mu_{\text{small}}=\mu_{\text{regular}}=\mu_{\text{aide}}$
(equal group means). Rows 1--3 have $p=.09,.26,.32$, all above .05, so equality is not rejected: the
pre-assignment traits are balanced, as random assignment predicts. Rows 4--6 are measured \emph{after}
assignment, so they are outcomes of the treatment: row 5 (15.10 against 22.40) shows the experiment
created the class-size gap (the ``first stage''), row 6 is the effect itself, and row 4 ($p=.02$) flags
differential attrition, a threat for later grades but not for kindergarten scores.}

\passo{(b) selection bias}{Write $Y_i=Y_{0i}+(Y_{1i}-Y_{0i})D_i$, so the observed difference in means is
\begin{equation}
E[Y_i\mid D_i=1]-E[Y_i\mid D_i=0]=E[Y_{1i}\mid D_i=1]-E[Y_{0i}\mid D_i=0].\label{q1:obs}
\end{equation}
Adding and subtracting $E[Y_{0i}\mid D_i=1]$ in \eqref{q1:obs} gives
\begin{equation}
E[Y_i\mid D_i=1]-E[Y_i\mid D_i=0]=\underbrace{E[Y_{1i}-Y_{0i}\mid D_i=1]}_{\text{effect on the treated}}
+\underbrace{E[Y_{0i}\mid D_i=1]-E[Y_{0i}\mid D_i=0]}_{\text{selection bias}}.\label{q1:dec}
\end{equation}
Random assignment makes $D_i$ independent of $(Y_{0i},Y_{1i})$:
\begin{equation}
E[Y_{0i}\mid D_i=1]=E[Y_{0i}\mid D_i=0],\qquad E[Y_{1i}-Y_{0i}\mid D_i=1]=E[Y_{1i}-Y_{0i}].\label{q1:ind}
\end{equation}
Substituting \eqref{q1:ind} into \eqref{q1:dec}, the selection-bias term is $(0)$ and the first term is
the average treatment effect. A regression of $Y_i$ on a treatment dummy returns exactly the difference in
group means, so $\hat\rho$ in column (1) estimates the average causal effect of a small class (relative to
a regular class with a part-time aide).}

\passo{(c) inference}{With $t=\hat\rho/\widehat{se}$ and the interval $\hat\rho\pm1.96\,\widehat{se}$,
column (1) gives
\begin{align}
t&=\frac{(4.82)}{(2.19)}=2.2009,\label{q1:t1}\\
\hat\rho\pm1.96\,\widehat{se}&=(4.82)\pm(1.96)(2.19)=4.82\pm4.2924=[0.5276,\;9.1124].\label{q1:ci1}
\end{align}
Column (4) gives
\begin{align}
t&=\frac{(5.37)}{(1.19)}=4.5126,\label{q1:t4}\\
\hat\rho\pm1.96\,\widehat{se}&=(5.37)\pm(1.96)(1.19)=5.37\pm2.3324=[3.0376,\;7.7024].\label{q1:ci4}
\end{align}
Both exceed 1.96, so the small-class effect is significant. The regular/aide coefficient has
$t=(.12)/(2.23)=0.0538$: no evidence that an aide matters.}

\passo{(d) covariates}{Let $X_i$ be the added covariates and $\pi$ the slope of $X_i$ on $D_i$. By the
omitted-variables formula (MHE \S3.2.2, pp.~44--47),
\begin{equation}
\rho^{\text{short}}=\rho^{\text{long}}+\gamma'\pi,\label{q1:ovb}
\end{equation}
where $\gamma$ is the effect of $X_i$ on $Y_i$. Randomization makes $\pi\approx0$, so by \eqref{q1:ovb}
the coefficient barely changes (5.37 against 5.36). The covariates do explain scores, which lowers the
residual variance. With homoskedastic errors,
\begin{equation}
\widehat{se}(\hat\rho)\approx\sqrt{\frac{(1-R^2)\,V(Y_i)}{n\,V(\tilde D_i)}},\label{q1:se}
\end{equation}
where $\tilde D_i$ is the residual of $D_i$ on the other regressors; since $D_i$ is unrelated to $X_i$,
$V(\tilde D_i)$ is unchanged. Dividing \eqref{q1:se} for column (3) by \eqref{q1:se} for column (2):
\begin{align}
\frac{\widehat{se}_3}{\widehat{se}_2}&\approx\sqrt{\frac{1-R^2_3}{1-R^2_2}}=\sqrt{\frac{1-(.31)}{1-(.25)}}=\sqrt{\frac{0.69}{0.75}}=0.95917,\label{q1:ratio}\\
\widehat{se}_3&\approx(1.26)(0.95917)=1.20855,\label{q1:se3}
\end{align}
which rounds to the reported 1.21. (The table's clustered standard errors make \eqref{q1:se} an approximation.)}

\passo{(e) effect size}{Setting $.2\sigma$ equal to the column (4) estimate,
\begin{equation}
\sigma=\frac{(5.37)}{(0.2)}=26.85.\label{q1:sigma}
\end{equation}
A percentile rank spread evenly over $[0,100]$ has standard deviation $100/\sqrt{12}=28.8675$, so 26.85 is
plausible. The raw gap in Table 2.2.1, row 6, is $(54.70)-(48.90)=5.80$, close to the regression estimates.}

\intuicao{Randomization removes selection bias, so the simple comparison is already causal; covariates
buy precision, not identification. If adding controls had moved $\hat\rho$ a lot, that would have been a
sign the randomization failed.}
\refbib{MHE \S2.1--2.3 (pp.~10--18), selection bias and regression analysis of experiments; \S3.2.2 (pp.~44--47), OVB formula.}
\end{sol}

% ---------------------------------------------------------------- Q2
\questao{Returns to schooling and omitted-variables bias (MHE Table 3.2.1).}
\fonte{Angrist and Pischke (2008), \textit{Mostly Harmless Econometrics}, \S3.2.2, Table 3.2.1, printed
p.~62. The source has no exercise; the statement is original. Table reproduced verbatim from the PDF.}

\begin{center}\small
\textit{Table 3.2.1: Estimates of the returns to education for men in the NLSY}\\[2pt]
\begin{tabular}{lccccc}
\toprule
 & (1) & (2) & (3) & (4) & (5)\\
\midrule
Controls: & None & Age & Col.~(2) and & Col.~(3) and & Col.~(4), with\\
 & & dummies & additional & AFQT score & occupation\\
 & & & controls* & & dummies\\
\midrule
 & 0.132 & 0.131 & 0.114 & 0.087 & 0.066\\
 & (0.007) & (0.007) & (0.007) & (0.009) & (0.010)\\
\bottomrule
\end{tabular}\\[2pt]
\parbox{0.9\linewidth}{\footnotesize Notes: ``Data are from the National Longitudinal Survey of Youth (1979
cohort, 2002 survey). The table reports the coefficient on years of schooling in a regression of log wages
on years of schooling and the indicated controls. Standard errors are shown in parentheses. The sample is
restricted to men and weighted by NLSY sampling weights. The sample size is 2434.
*Additional controls are mother's and father's years of schooling and dummy variables for race and Census region.''}
\end{center}

\begin{enumerate}[label=(\alph*)]
\item Write the long regression of log wages on schooling, the column (3) controls and AFQT, and the
auxiliary regression of AFQT on schooling and the same controls. Derive the omitted-variables-bias
(OVB) formula that links the schooling coefficients in columns (3) and (4).
\item Compute the change in the schooling coefficient attributable to AFQT, and to the family-background
controls (column (2) to column (3)). What do the signs say about the correlations involved?
\item Suppose the auxiliary regression of AFQT (in standard-deviation units) on schooling, with the column
(3) controls, had slope 0.15 (a hypothetical value, not in the book). What effect of one standard deviation
of AFQT on log wages does the OVB formula then imply?
\item Column (5) is half of column (1). Why is column (5) \emph{not} the most credible estimate of the
return to schooling?
\item Give a 95\% confidence interval for the return to schooling in column (4).
\end{enumerate}
\resolespaco
\begin{sol}
\begin{notacao}
\item $Y_i$: log wage; $S_i$: years of schooling; $A_i$: AFQT score (standard-deviation units);
$X_i$: the column (3) controls.
\item $\rho$: schooling coefficient in the long regression (with $A_i$); $\rho^s$: in the short one
(without $A_i$); $\gamma$: AFQT coefficient; $\delta_{AS}$: slope of $A_i$ on $S_i$ given $X_i$.
\end{notacao}

\passo{(a) OVB formula}{The long, auxiliary and short regressions are
\begin{align}
Y_i&=\alpha+\rho S_i+X_i'\beta+\gamma A_i+e_i,\label{q2:long}\\
A_i&=\delta_0+\delta_{AS}S_i+X_i'\pi+v_i,\label{q2:aux}\\
Y_i&=\alpha^s+\rho^sS_i+X_i'\beta^s+e^s_i.\label{q2:short}
\end{align}
Substituting \eqref{q2:aux} into \eqref{q2:long}:
\begin{equation}
Y_i=\alpha+\rho S_i+X_i'\beta+\gamma\left(\delta_0+\delta_{AS}S_i+X_i'\pi+v_i\right)+e_i.\label{q2:sub}
\end{equation}
Collecting terms in \eqref{q2:sub}:
\begin{equation}
Y_i=(\alpha+\gamma\delta_0)+(\rho+\gamma\delta_{AS})S_i+X_i'(\beta+\gamma\pi)+(\gamma v_i+e_i).\label{q2:coll}
\end{equation}
The error $\gamma v_i+e_i$ is uncorrelated with $S_i$ and $X_i$ ($v_i$ is an OLS residual, $e_i$ is the
long-regression residual), so \eqref{q2:coll} is the short regression \eqref{q2:short} and
\begin{equation}
\rho^s=\rho+\gamma\,\delta_{AS}.\label{q2:ovb}
\end{equation}}

\passo{(b) the numbers}{Column (3) is the short regression and column (4) the long one. From \eqref{q2:ovb}:
\begin{equation}
\gamma\,\delta_{AS}=\rho^s-\rho=(0.114)-(0.087)=0.027.\label{q2:biasA}
\end{equation}
The bias is positive, so $\gamma$ and $\delta_{AS}$ have the same sign: AFQT is positively related to both
wages and schooling. The family-background block lowers the coefficient by $(0.131)-(0.114)=0.017$, again
positive: parents' schooling raises both the son's schooling and his wage.}

\passo{(c) implied AFQT effect}{Solving \eqref{q2:biasA} for $\gamma$ with $\delta_{AS}=0.15$:
\begin{equation}
\gamma=\frac{0.027}{\delta_{AS}}=\frac{0.027}{(0.15)}=0.18,\label{q2:gamma}
\end{equation}
i.e.\ one standard deviation of AFQT would raise wages by about 18 log points, holding schooling fixed.}

\passo{(d) bad control}{Occupation is itself an outcome of schooling. Conditioning on it compares men
within occupations, which (i) removes a channel through which schooling raises wages and (ii) creates
selection bias, because the men with low schooling who reach a given occupation differ in unobservables
from those with high schooling in the same occupation. Column (5), $0.066/0.132=0.50$ of column (1), mixes
these two effects and has no causal interpretation. Column (4) is preferred, keeping in mind that AFQT may
itself be partly affected by schooling (a ``proxy control'').}

\passo{(e) interval}{Column (4): $(0.087)\pm(1.96)(0.009)=0.087\pm0.01764=[0.06936,\;0.10464]$.}

\intuicao{Each block of controls lowers the coefficient because each omitted factor is positively
correlated with both schooling and wages. The OVB formula tells you the direction of the bias before you
run anything; the bad-control argument tells you where to stop adding controls.}
\refbib{MHE \S3.2.2 (pp.~44--47), OVB formula; \S3.2.3 (pp.~47--50), bad control; W \S3.3, omitted-variable bias.}
\end{sol}

% ---------------------------------------------------------------- Q3
\questao{Wald and 2SLS estimates: quarter of birth and the draft lottery (MHE Tables 4.1.1--4.1.3).}
\fonte{Angrist and Pischke (2008), \textit{Mostly Harmless Econometrics}, \S4.1, Tables 4.1.1, 4.1.2 and
4.1.3, printed pp.~92, 96 and 97. The source has no exercise; the statement is original. Table numbers
reproduced verbatim from the PDF (notes shortened).}

\begin{center}\small
\textit{Table 4.1.1: 2SLS estimates of the economic returns to schooling}\\[2pt]
\resizebox{\linewidth}{!}{%
\begin{tabular}{lcccccccc}
\toprule
 & \multicolumn{2}{c}{OLS} & \multicolumn{6}{c}{2SLS}\\
\cmidrule(lr){2-3}\cmidrule(lr){4-9}
 & (1) & (2) & (3) & (4) & (5) & (6) & (7) & (8)\\
\midrule
Years of education & 0.075 & 0.072 & 0.103 & 0.112 & 0.106 & 0.108 & 0.089 & 0.061\\
 & (0.0004) & (0.0004) & (0.024) & (0.021) & (0.026) & (0.019) & (0.016) & (0.031)\\[3pt]
\textit{Covariates:} & & & & & & & & \\
Age (in quarters) & & & & & & & & $\checkmark$\\
Age (in quarters) squared & & & & & & & & $\checkmark$\\
9 year of birth dummies & & $\checkmark$ & & & $\checkmark$ & $\checkmark$ & $\checkmark$ & $\checkmark$\\
50 state of birth dummies & & $\checkmark$ & & & $\checkmark$ & $\checkmark$ & $\checkmark$ & $\checkmark$\\[3pt]
\textit{Instruments:} & & & \shortstack{dummy\\for\\QOB=1} & \shortstack{dummy\\for\\QOB=1\\or\\QOB=2}
 & \shortstack{dummy\\for\\QOB=1} & \shortstack{full set\\of QOB\\dummies}
 & \shortstack{full set\\of QOB\\dummies\\int. with\\year of\\birth\\dummies}
 & \shortstack{full set\\of QOB\\dummies\\int. with\\year of\\birth\\dummies}\\
\bottomrule
\end{tabular}}\\[2pt]
\parbox{0.92\linewidth}{\footnotesize Notes (excerpt): ``The table reports OLS and 2SLS estimates of the returns to
schooling using the the Angrist and Krueger (1991) 1980 Census sample. [\dots] The sample size is 329,509.
Robust standard errors are reported in parentheses.''}

\medskip
\textit{Table 4.1.2: Wald estimates of the returns to schooling using quarter of birth instruments}\\[2pt]
\begin{tabular}{lccc}
\toprule
 & (1) & (2) & (3)\\
 & Born in the 1st & Born in the 3rd & Difference\\
 & or 2nd quarter of & or 4th quarter of & (std.\ error)\\
 & year & year & (1)-(2)\\
\midrule
ln (weekly wage) & 5.8916 & 5.9051 & -0.01349\\
 & & & (0.00337)\\
Years of education & 12.6881 & 12.8394 & -0.1514\\
 & & & (0.0162)\\
Wald estimate of return to education & & & 0.0891\\
 & & & (0.0210)\\
OLS estimate of return to education & & & 0.0703\\
 & & & (0.0005)\\
\bottomrule
\end{tabular}\\[2pt]
\parbox{0.92\linewidth}{\footnotesize Notes (excerpt): ``Adapted from a re-analysis of Angrist and Krueger (1991) by
Angrist and Imbens (1995). [\dots] The sample size is 329,509.''}

\medskip
\textit{Table 4.1.3: Wald estimates of the effects of military service on the earnings of white men born in 1950}\\[2pt]
\begin{tabular}{lccccc}
\toprule
 & \multicolumn{2}{c}{Earnings} & \multicolumn{2}{c}{Veteran Status} & Wald Estimate\\
\cmidrule(lr){2-3}\cmidrule(lr){4-5}
Earnings year & Mean & Eligibility Effect & Mean & Eligibility Effect & of Veteran Effect\\
 & (1) & (2) & (3) & (4) & (5)\\
\midrule
1981 & 16,461 & -435.8 & 0.267 & 0.159 & -2,741\\
 & & (210.5) & & (0.040) & (1,324)\\
1971 & 3,338 & -325.9 & & & -2050\\
 & & (46.6) & & & (293)\\
1969 & 2,299 & -2.0 & & & \\
 & & (34.5) & & & \\
\bottomrule
\end{tabular}\\[2pt]
\parbox{0.92\linewidth}{\footnotesize Notes (excerpt): ``Adapted from Angrist (1990), Tables 2 and 3. Standard errors are
shown in parentheses. [\dots] Figures are in nominal dollars. [\dots] There are about 13,500 individuals in the sample.''}
\end{center}

\begin{enumerate}[label=(\alph*)]
\item Derive the Wald estimator for a binary instrument and state the assumptions under which it
identifies a causal effect.
\item Table 4.1.2: compute the Wald estimate from column (3) and check it against 0.0891. Approximate its
standard error treating the first stage as known, and compare with 0.0210.
\item Table 4.1.3: compute the Wald estimates for 1981 and 1971 from columns (2) and (4), their approximate
standard errors, and each as a share of mean earnings. What is the 1969 row for?
\item Table 4.1.1: is the 2SLS estimate in column (3) significantly different from the OLS estimate in
column (1)? Use a Hausman-type statistic with $V(\hat\rho_{IV}-\hat\rho_{OLS})=V(\hat\rho_{IV})-V(\hat\rho_{OLS})$.
Why is the standard error in column (8) larger than in columns (6)--(7), even though (8) uses as many instruments?
\end{enumerate}
\resolespaco
\begin{sol}
\begin{notacao}
\item $Y_i$: outcome (log weekly wage, or annual earnings in dollars); $D_i$: treatment (years of
schooling, or veteran status $\in\{0,1\}$); $Z_i\in\{0,1\}$: instrument (born in quarters 1--2, or draft-eligible).
\item $\rho$: causal effect of $D_i$; $\eta_i$: structural error; $\pi=E[D_i\mid Z_i=1]-E[D_i\mid Z_i=0]$:
first stage; $\phi=E[Y_i\mid Z_i=1]-E[Y_i\mid Z_i=0]$: reduced form.
\end{notacao}

\passo{(a) Wald estimator}{Take the constant-effect model and the instrument assumption
\begin{equation}
Y_i=\alpha+\rho D_i+\eta_i,\qquad E[\eta_i\mid Z_i]=E[\eta_i].\label{q3:struct}
\end{equation}
Conditioning \eqref{q3:struct} on $Z_i=z$:
\begin{equation}
E[Y_i\mid Z_i=z]=\alpha+\rho\,E[D_i\mid Z_i=z]+E[\eta_i].\label{q3:cond}
\end{equation}
Subtracting \eqref{q3:cond} at $z=0$ from \eqref{q3:cond} at $z=1$:
\begin{equation}
\phi=\rho\,\pi.\label{q3:rf}
\end{equation}
Dividing \eqref{q3:rf} by $\pi\neq0$:
\begin{equation}
\rho=\frac{\phi}{\pi}=\frac{E[Y_i\mid Z_i=1]-E[Y_i\mid Z_i=0]}{E[D_i\mid Z_i=1]-E[D_i\mid Z_i=0]}.\label{q3:wald}
\end{equation}
Assumptions: $Z_i$ as good as randomly assigned (independence), $Z_i$ affects $Y_i$ only through $D_i$
(exclusion), and $\pi\neq0$ (first stage). With heterogeneous effects and monotonicity (no defiers),
\eqref{q3:wald} is the local average treatment effect on compliers (MHE \S4.4).}

\passo{(b) quarter of birth}{Substituting column (3) of Table 4.1.2 into \eqref{q3:wald}:
\begin{equation}
\hat\rho=\frac{(-0.01349)}{(-0.1514)}=0.089102,\label{q3:wqob}
\end{equation}
which rounds to the reported 0.0891. (The rounded means give $5.8916-5.9051=-0.0135$ and
$12.6881-12.8394=-0.1513$.) Treating $\pi$ as known, $\widehat{se}(\hat\rho)\approx\widehat{se}(\hat\phi)/|\hat\pi|$:
\begin{equation}
\widehat{se}(\hat\rho)\approx\frac{(0.00337)}{(0.1514)}=0.022259,\label{q3:seqob}
\end{equation}
close to 0.0210; the delta method adds the sampling noise in $\hat\pi$ and its covariance with $\hat\phi$.
The Wald estimate (0.0891) exceeds OLS (0.0703).}

\passo{(c) draft lottery}{Substituting columns (2) and (4) of Table 4.1.3 into \eqref{q3:wald}:
\begin{align}
\hat\rho_{1981}&=\frac{(-435.8)}{(0.159)}=-2740.88,\label{q3:w81}\\
\hat\rho_{1971}&=\frac{(-325.9)}{(0.159)}=-2049.69,\label{q3:w71}
\end{align}
matching $-2{,}741$ and $-2050$. Dividing the reduced-form standard errors by the first stage, as in
\eqref{q3:seqob}: $(210.5)/(0.159)=1323.90$ and $(46.6)/(0.159)=293.08$, matching (1,324) and (293). As
shares of mean earnings: $2740.88/16461=0.16651$ in 1981 and $2049.69/3338=0.61405$ in 1971, so service
cut earnings by about 17\% ten years later and by much more while the men were still in the army. The
1969 row is a placebo: earnings before anyone born in 1950 was drafted cannot depend on veteran status,
and indeed $t=(-2.0)/(34.5)=-0.0580$. A zero there supports the exclusion restriction.}

\passo{(d) OLS against 2SLS}{The Hausman-type statistic with the given variance:
\begin{equation}
H=\frac{(0.103)-(0.075)}{\sqrt{(0.024)^2-(0.0004)^2}}=\frac{0.028}{0.023997}=1.1668<1.96,\label{q3:haus}
\end{equation}
so the difference is not significant: quarter of birth is too weak a source of variation to separate the two.
Column (8) adds age in quarters and its square. Within a year of birth, age in quarters is a linear
function of quarter of birth, so these controls absorb much of the instrument's variation: the first stage
weakens and the standard error rises (0.031 against 0.016--0.019). Columns (7)--(8) also use 30 instruments,
which pulls 2SLS towards OLS when the instruments are weak (the many-weak-instruments bias discussed in MHE Ch.~4).}

\intuicao{The Wald estimator rescales the effect of the instrument on the outcome by its effect on the
treatment. The QOB first stage is tiny (0.15 years), so small reduced-form effects translate into
imprecise IV estimates; the draft first stage is larger (16 points), so the estimates are informative.}
\refbib{MHE \S4.1 (pp.~84--102), IV and the Wald estimator; \S4.4 (pp.~111--129), LATE; W \S15.1 (pp.~462--471).}
\end{sol}

% ---------------------------------------------------------------- Q4
\questao{Card and Krueger: difference-in-differences by hand (MHE Table 5.2.1).}
\fonte{Angrist and Pischke (2008), \textit{Mostly Harmless Econometrics}, \S5.2, Table 5.2.1, printed
p.~171. The source has no exercise; the statement is original. Table reproduced verbatim from the PDF.}

\begin{center}\small
\textit{Table 5.2.1: Average employment per store before and after the New Jersey minimum wage increase}\\[2pt]
\begin{tabular}{llccc}
\toprule
 & & PA & NJ & Difference, NJ-PA\\
 & Variable & (i) & (ii) & (iii)\\
\midrule
1. & FTE employment before, & 23.33 & 20.44 & -2.89\\
 & all available observations & (1.35) & (0.51) & (1.44)\\
2. & FTE employment after, & 21.17 & 21.03 & -0.14\\
 & all available observations & (0.94) & (0.52) & (1.07)\\
3. & Change in mean FTE & -2.16 & 0.59 & 2.76\\
 & employment & (1.25) & (0.54) & (1.36)\\
\bottomrule
\end{tabular}\\[2pt]
\parbox{0.9\linewidth}{\footnotesize Notes (excerpt): ``Adapted from Card and Krueger (1994), Table 3. The table
reports average full-time equivalent (FTE) employment at restaurants in Pennsylvania and New Jersey before
and after a minimum wage increase in New Jersey. [\dots] Standard errors are reported in parentheses''}
\end{center}

\begin{enumerate}[label=(\alph*)]
\item Compute the difference-in-differences estimate twice: as the difference of the two states' changes,
and as the change in the NJ--PA gap. Reconcile with 2.76.
\item Write the regression $Y_{ist}=\alpha+\gamma NJ_s+\lambda d_t+\delta(NJ_s\cdot d_t)+\varepsilon_{ist}$ and
give the value of each coefficient implied by the table. State the identifying assumption.
\item Compute the $t$ statistic and a 95\% interval for the DiD estimate, and show that its standard
error is what independent samples in the two states imply.
\item Is the sign of the estimate what a competitive labor-demand model predicts?
\end{enumerate}
\resolespaco
\begin{sol}
\begin{notacao}
\item $Y_{ist}$: FTE employment of store $i$ in state $s\in\{PA,NJ\}$ at $t\in\{\text{before},\text{after}\}$;
$\bar Y_{st}$: cell mean; $Y_{0ist}$: employment without the minimum-wage rise.
\item $NJ_s$, $d_t$: dummies for New Jersey and for the after period; $\gamma_s,\lambda_t$: state and time
effects; $\delta$: causal effect of the minimum wage on FTE employment.
\end{notacao}

\passo{(a) two ways}{The DiD estimator is
\begin{equation}
\hat\delta=(\bar Y_{NJ,a}-\bar Y_{NJ,b})-(\bar Y_{PA,a}-\bar Y_{PA,b}).\label{q4:did}
\end{equation}
Substituting rows 1--2:
\begin{align}
\hat\delta&=\big((21.03)-(20.44)\big)-\big((21.17)-(23.33)\big)\label{q4:sub1}\\
&=(0.59)-(-2.16)=2.75.\label{q4:res1}
\end{align}
Rearranging \eqref{q4:did} as the change in the gap and substituting column (iii):
\begin{equation}
\hat\delta=(\bar Y_{NJ,a}-\bar Y_{PA,a})-(\bar Y_{NJ,b}-\bar Y_{PA,b})=(-0.14)-(-2.89)=2.75.\label{q4:res2}
\end{equation}
The table's 2.76 comes from unrounded cell means.}

\passo{(b) regression}{Common trends means
\begin{equation}
E[Y_{0ist}\mid s,t]=\gamma_s+\lambda_t.\label{q4:ct}
\end{equation}
With $Y_{ist}=Y_{0ist}+\delta\,NJ_s d_t$, the NJ change is $(\lambda_a-\lambda_b)+\delta$ and the PA change
is $\lambda_a-\lambda_b$; their difference is $\delta$, which is what \eqref{q4:did} estimates. In the
regression the four cells are $\alpha$, $\alpha+\gamma$, $\alpha+\lambda$, $\alpha+\gamma+\lambda+\delta$, so
\begin{equation}
\hat\alpha=23.33,\quad\hat\gamma=(20.44)-(23.33)=-2.89,\quad\hat\lambda=(21.17)-(23.33)=-2.16,\quad\hat\delta=2.75.\label{q4:coefs}
\end{equation}
The identifying assumption is \eqref{q4:ct}: without the rise, NJ employment would have moved like PA employment.}

\passo{(c) inference}{If the two states' samples are independent, the variance of a difference is the sum of the variances:
\begin{equation}
\widehat{se}(\hat\delta)=\sqrt{(1.25)^2+(0.54)^2}=\sqrt{1.5625+0.2916}=\sqrt{1.8541}=1.36165,\label{q4:se}
\end{equation}
the reported 1.36. Then $t=(2.76)/(1.36)=2.0294$ and $2.76\pm(1.96)(1.36)=2.76\pm2.6656=[0.0944,\;5.4256]$:
significant at 5\%, but the interval is wide.}

\passo{(d) interpretation}{A competitive model predicts lower employment in NJ, so $\delta<0$. The estimate
is positive: employment fell in PA and rose slightly in NJ. Either the market is not competitive
(monopsony), or common trends fails; MHE shows that longer payroll data from the same states weaken the
original result, which is why pre-period trends should be checked.}

\intuicao{DiD uses the control state's change as the counterfactual change for the treated state; the
state effect $\gamma$ absorbs any permanent NJ--PA gap, so only differential trends can bias $\hat\delta$.}
\refbib{MHE \S5.2 (pp.~169--181), differences-in-differences; W \S13.2 (pp.~406--412).}
\end{sol}

% ---------------------------------------------------------------- Q5
\questao{Incumbency and the sharp regression-discontinuity design (MHE Figure 6.1.2).}
\fonte{Angrist and Pischke (2008), \textit{Mostly Harmless Econometrics}, \S6.1, Figure 6.1.2 and text,
printed pp.~194--195. The source has no exercise; the statement is original. Caption and numbers from the
text reproduced verbatim; the two fitted values at the cutoff are read off the figure (not printed in the book).}

\smallskip
\noindent Figure 6.1.2 (caption, verbatim): ``Probability of winning an election by past and future vote
share (from Lee, 2008). (a) Candidate's probability of winning election $t+1$, by margin of victory in
election $t$: local averages and parametric fit. (b) Candidate's accumulated number of past election
victories, by margin of victory in election $t$: local averages and parametric fit.'' The text states
that the winner is determined by $D_i=1(x_i\geq.0)$, where $x_i$ is the vote share margin of victory; that
the dots are ``the average win rate in non-overlapping windows of share margins that are .005 wide''; that
the lines are fitted from ``a Logit model for the probability of winning as a function of the cutoff
indicator $D_i=1(x_i\geq0)$, a $4^{th}$-order polynomial in $x_i$, and interactions between the polynomial
terms and $D_i$''; and that ``incumbency appears to raise party re-election probabilities by about 40
percentage points.'' In panel (a) the fitted curve reaches about 0.14 just left of zero and about 0.58
just right of zero; in panel (b) the two fitted curves nearly meet at zero (both close to one past victory).

\begin{enumerate}[label=(\alph*)]
\item Write the sharp RD model with a polynomial control and define the estimand as a jump in a
conditional expectation. State the identifying assumption.
\item Estimate the incumbency effect from panel (a) and compare it with ``about 40 percentage points''.
\item What does panel (b) test, and what does it show?
\item Give two reasons why an RD estimate could be spurious here, and a check for each.
\end{enumerate}
\resolespaco
\begin{sol}
\begin{notacao}
\item $x_i$: Democratic vote-share margin in election $t$ (running variable, cutoff $0$);
$D_i=1(x_i\geq0)$: Democrat won election $t$ (incumbency).
\item $Y_i$: Democrat wins election $t+1$; $Y_{0i},Y_{1i}$: potential outcomes; $f(\cdot)$: smooth
function of $x_i$ (a polynomial); $\rho$: incumbency effect at the cutoff.
\end{notacao}

\passo{(a) model}{Assume $E[Y_{0i}\mid x_i]=f(x_i)$ and a constant effect $\rho$:
\begin{equation}
Y_i=f(x_i)+\rho D_i+\eta_i,\qquad D_i=1(x_i\geq0).\label{q5:model}
\end{equation}
Taking limits of $E[Y_i\mid x_i]$ in \eqref{q5:model} from each side of $0$, with $f$ continuous:
\begin{align}
\lim_{x\downarrow0}E[Y_i\mid x_i=x]&=f(0)+\rho,\label{q5:right}\\
\lim_{x\uparrow0}E[Y_i\mid x_i=x]&=f(0).\label{q5:left}
\end{align}
Subtracting \eqref{q5:left} from \eqref{q5:right}:
\begin{equation}
\rho=\lim_{x\downarrow0}E[Y_i\mid x_i=x]-\lim_{x\uparrow0}E[Y_i\mid x_i=x].\label{q5:jump}
\end{equation}
Identifying assumption: $E[Y_{0i}\mid x_i]$ and $E[Y_{1i}\mid x_i]$ are continuous at $0$. Because $D_i$ is
a deterministic function of $x_i$, there is no other confounder once $f(x_i)$ is controlled for.}

\passo{(b) the jump}{Substituting the fitted values at the cutoff into \eqref{q5:jump}:
\begin{equation}
\hat\rho\approx(0.58)-(0.14)=0.44,\label{q5:est}
\end{equation}
an incumbency advantage of roughly 40 percentage points, as the text says.}

\passo{(c) panel (b)}{Past victories are determined before election $t$, so they cannot be caused by
$D_i$. If close elections are as good as random, their conditional mean must be continuous at $0$:
\eqref{q5:jump} applied to this variable should be zero. The two fitted curves nearly meet at $0$, so there
is no visible jump: this balance check supports the design, as a covariate-balance table does in an experiment.}

\passo{(d) threats}{(i) Misspecified $f$: a nonlinear $E[Y_{0i}\mid x_i]$ fitted with too rigid a polynomial
can show up as a jump. Check: higher-order and side-specific polynomials (the interactions with $D_i$ in
the footnote), and local estimates in narrow windows around $0$. (ii) Manipulation: if candidates could
sort just above $0$ (recounts, fraud), units on the two sides would differ. Check: the density of $x_i$
near $0$ (McCrary test) and balance of pre-determined covariates such as panel (b).}

\intuicao{Bare winners and bare losers are alike in everything but the election outcome, so the
discontinuity at a zero margin isolates what winning itself does to the next election.}
\refbib{MHE \S6.1 (pp.~189--196), sharp RD; class slides and site example on RDD (minimum drinking age and motor-vehicle deaths).}
\end{sol}

% ---------------------------------------------------------------- Q6
\questao{Maimonides' rule and fuzzy RD as IV (MHE Table 6.2.1).}
\fonte{Angrist and Pischke (2008), \textit{Mostly Harmless Econometrics}, \S6.2, Table 6.2.1 and text,
printed pp.~199--202. The source has no exercise; the statement is original. Table and the rule
reproduced verbatim from the PDF.}

\smallskip
\noindent The text defines the predicted class size from Maimonides' rule (class size capped at 40) as
\[
m_{sc}=\frac{e_s}{\operatorname{int}\left[\frac{(e_s-1)}{40}\right]+1},
\]
``where int(x) is the integer part of a real number, x'', $e_s$ is enrollment in the grade, and
estimates by 2SLS ``an equation like $y_{isc}=\alpha_0+\alpha_1d_s+\beta_1e_s+\beta_2e_s^2+\dots+\beta_pe_s^p+\rho n_{sc}+\eta_{isc}$''
(6.2.6), with $n_{sc}$ class size, $d_s$ the proportion of disadvantaged students, and $m_{sc}$ as instrument.

\begin{center}\small
\textit{Table 6.2.1: OLS and fuzzy RD estimates of the effects of class size on fifth grade math scores}\\[2pt]
\resizebox{\linewidth}{!}{%
\begin{tabular}{lcccccccc}
\toprule
 & \multicolumn{3}{c}{OLS} & \multicolumn{5}{c}{2SLS}\\
\cmidrule(lr){2-4}\cmidrule(lr){5-9}
 & & & & \multicolumn{2}{c}{Full sample} & \multicolumn{3}{c}{Discontinuity samples}\\
\cmidrule(lr){5-6}\cmidrule(lr){7-9}
 & & & & & & \multicolumn{2}{c}{+/- 5} & +/- 3\\
 & (1) & (2) & (3) & (4) & (5) & (6) & (7) & (8)\\
\midrule
\textit{Mean score} & & 67.3 & & 67.3 & & 67.0 & & 67.0\\
\textit{(s.d.)} & & (9.6) & & (9.6) & & (10.2) & & (10.6)\\
\textit{Regressors} & & & & & & & & \\
Class size & .322 & .076 & .019 & -.230 & -.261 & -.185 & -.443 & -.270\\
 & (.039) & (.036) & (.044) & (.092) & (.113) & (.151) & (.236) & (.281)\\
Percent disadvantaged & & -.340 & -.332 & -.350 & -.350 & -.459 & -.435 & \\
 & & (.018) & (.018) & (.019) & (.019) & (.049) & (.049) & \\
Enrollment & & & .017 & .041 & .062 & & .079 & \\
 & & & (.009) & (.012) & (.037 & & (.036) & \\
Enrollment squared/100 & & & & & -.010 & & & \\
 & & & & & (.016) & & & \\
Segment 1 (enrollment 36-45) & & & & & & & & -12.6\\
 & & & & & & & & (3.80)\\
Segment 2 (enrollment 76-85) & & & & & & & & -2.89\\
 & & & & & & & & (2.41)\\
Root MSE & 9.36 & 8.32 & 8.30 & 8.40 & 8.42 & 8.79 & 9.10 & 10.2\\
R-squared & .048 & .249 & .252 & & & & & \\
N & & 2,018 & & & 2,018 & & 471 & 302\\
\bottomrule
\end{tabular}}\\[2pt]
\parbox{0.92\linewidth}{\footnotesize Notes: ``Adapted from Angrist and Lavy (1999). The table reports estimates of
equation (6.2.6) in the text using class averages. Standard errors, reported in parentheses, are corrected
for within-school correlation.'' (The missing parenthesis in ``(.037'' is in the original.)}
\end{center}

\begin{enumerate}[label=(\alph*)]
\item Compute $m_{sc}$ for enrollments $e_s=40,41,80,81,120,121$. Why does the rule generate a
\emph{fuzzy} rather than a sharp RD?
\item Write the first and second stages of the 2SLS model. Which variables play the roles of the
instrument, the running variable and the treatment?
\item Explain the sign and the movement of the OLS class-size coefficient across columns (1)--(3).
\item Using column (4) and the text's round slope of $-.25$, how much would a 7-pupil cut in class size
raise math scores, in points and in standard deviations of the class-average score (9.6)? Give the $t$
statistics in columns (4) and (8).
\item Compare the ratio of the standard errors in columns (8) and (5) with $\sqrt{2018/302}$. What does the
comparison say?
\end{enumerate}
\resolespaco
\begin{sol}
\begin{notacao}
\item $e_s$: grade enrollment in school $s$ (pupils, integer $\geq1$); $m_{sc}$: class size predicted by the
rule; $n_{sc}$: actual class size; $d_s$: percent disadvantaged.
\item $y_{sc}$: class-average math score; $\rho$: effect of one more pupil per class on the score;
$\pi_m$: first-stage coefficient on $m_{sc}$; $\operatorname{int}(x)$: integer part of $x$.
\end{notacao}

\passo{(a) the rule}{Substituting each enrollment into the rule:
\begin{align}
m(40)&=\frac{40}{\operatorname{int}(39/40)+1}=\frac{40}{(0)+1}=40,&
m(41)&=\frac{41}{\operatorname{int}(40/40)+1}=\frac{41}{(1)+1}=20.5,\label{q6:m1}\\
m(80)&=\frac{80}{\operatorname{int}(79/40)+1}=\frac{80}{(1)+1}=40,&
m(81)&=\frac{81}{\operatorname{int}(80/40)+1}=\frac{81}{(2)+1}=27,\label{q6:m2}\\
m(120)&=\frac{120}{\operatorname{int}(119/40)+1}=\frac{120}{(2)+1}=40,&
m(121)&=\frac{121}{\operatorname{int}(120/40)+1}=\frac{121}{(3)+1}=30.25.\label{q6:m3}
\end{align}
Predicted class size drops sharply after each multiple of 40. The design is fuzzy because actual class size
$n_{sc}$ does not follow the rule exactly (some schools split grades below 40), so crossing a cutoff changes
the \emph{expected} class size, not class size itself with certainty.}

\passo{(b) 2SLS}{The first stage and the second stage are
\begin{align}
n_{sc}&=\pi_0+\pi_1d_s+\pi_2e_s+\dots+\pi_m\,m_{sc}+\xi_{sc},\label{q6:fs}\\
y_{sc}&=\alpha_0+\alpha_1d_s+\beta_1e_s+\dots+\rho\,n_{sc}+\eta_{sc}.\label{q6:ss}
\end{align}
The instrument is $m_{sc}$ (the role of $t_i$), the running variable is $e_s$ (the role of $x_i$, controlled
smoothly), and the treatment is $n_{sc}$ (the role of $d_i$). Identification comes from the jumps in
$m_{sc}$ at 40, 80, 120 conditional on smooth functions of $e_s$.}

\passo{(c) OLS}{Column (1), $+.322$: larger classes are found in larger, more advantaged schools, and
weak pupils are often put in smaller classes, so class size is positively correlated with omitted
determinants of scores. Controlling for percent disadvantaged removes most of this (.076), and adding
enrollment brings it to .019, insignificant. By the OVB logic of Question 2, the OLS coefficients are
biased upwards.}

\passo{(d) magnitudes}{The effect of a change $\Delta n=-7$ is $\rho\,\Delta n$. With column (4) and with the
text's slope:
\begin{align}
\Delta y&=(-0.230)(-7)=1.61,\label{q6:d4}\\
\Delta y&=(-0.25)(-7)=1.75,\label{q6:d25}\\
\frac{\Delta y}{\text{s.d.}}&=\frac{(1.75)}{(9.6)}=0.18229,\label{q6:es}
\end{align}
about $.18\sigma$, close to the STAR estimate of about $.2\sigma$ in Question 1. The $t$ statistics in
columns (4) and (8) are
\begin{equation}
t_4=\frac{(-.230)}{(.092)}=-2.5,\qquad t_8=\frac{(-.270)}{(.281)}=-0.9609.\label{q6:t}
\end{equation}}

\passo{(e) precision}{The ratio of standard errors and the square root of the ratio of sample sizes are
\begin{align}
\frac{\widehat{se}_8}{\widehat{se}_5}&=\frac{(.281)}{(.113)}=2.4867,\label{q6:ser}\\
\sqrt{\frac{2018}{302}}&=\sqrt{6.6821}=2.5850.\label{q6:nr}
\end{align}
Standard errors scale with $1/\sqrt{N}$, so shrinking the sample to the $\pm3$ window costs precision
almost exactly in proportion to $\sqrt{N}$: the estimate stays near $-.25$ but becomes insignificant.}

\intuicao{The rule is an accident of administration that shifts class size up and down independently of
the school's quality once enrollment is controlled for smoothly; 2SLS rescales the score jumps at the
cutoffs by the class-size jumps, exactly like a Wald estimator.}
\refbib{MHE \S6.2 (pp.~196--202), fuzzy RD is IV; Question 3 for the Wald logic.}
\end{sol}

% =============================================================================
\bloco{2}{Instrumental variables with the CARD data}{Wooldridge Ch.~15 \,\textbullet\, Questions 7--8}

% ---------------------------------------------------------------- Q7
\questao{Wooldridge, Computer Exercise 15.C3.}
\fonte{Wooldridge (2016), \textit{Introductory Econometrics: A Modern Approach}, 6th ed., Chapter 15,
Computer Exercise C3, printed pp.~492--493, reproduced verbatim.}

\noindent\textbf{C3} Use the data in CARD for this exercise.
\begin{enumerate}[label=(\roman*)]
\item The equation we estimated in Example 15.4 can be written as
\[\log(wage)=\beta_0+\beta_1educ+\beta_2exper+\dots+u,\]
where the other explanatory variables are listed in Table 15.1. In order for IV to be consistent, the IV
for $educ$, $nearc4$, must be uncorrelated with $u$. Could $nearc4$ be correlated with things in the error
term, such as unobserved ability? Explain.
\item For a subsample of the men in the data set, an IQ score is available. Regress $IQ$ on $nearc4$ to
check whether average IQ scores vary by whether the man grew up near a four-year college. What do you conclude?
\item Now, regress $IQ$ on $nearc4$, $smsa66$, and the 1966 regional dummy variables $reg662$, \dots,
$reg669$. Are $IQ$ and $nearc4$ related after the geographic dummy variables have been partialled out?
Reconcile this with your findings from part (ii).
\item From parts (ii) and (iii), what do you conclude about the importance of controlling for $smsa66$ and
the 1966 regional dummies in the $\log(wage)$ equation?
\end{enumerate}
\resolespaco
\begin{sol}
\begin{notacao}
\item $IQ_i$: IQ score (available for 2,061 of the 3,010 men); $nearc4_i\in\{0,1\}$: grew up near a
four-year college; $smsa66_i$: lived in an SMSA in 1966; $reg66k_i$: 1966 region dummies ($k=2,\dots,9$).
\item $u_i$: error of the wage equation; $\theta$: coefficient on $nearc4$; $G_i$: the geographic dummies.
\end{notacao}

\noindent\textit{Code} (data \texttt{Aula/class-05-06/materials/card.csv}, CR line endings, ``.'' = missing):
\begingroup\footnotesize
\begin{verbatim}
import pandas as pd, statsmodels.formula.api as smf
card = pd.read_csv("Aula/class-05-06/materials/card.csv", lineterminator="\r", na_values=".")
iq = card.dropna(subset=["IQ"])                                   # n = 2,061
regs = "+".join(f"reg66{i}" for i in range(2, 10))
a = smf.ols("IQ ~ nearc4", iq).fit()                              # part (ii)
b = smf.ols("IQ ~ nearc4 + smsa66 + " + regs, iq).fit()           # part (iii)
print(b.f_test("smsa66=0," + ",".join(f"reg66{i}=0" for i in range(2, 10))))
\end{verbatim}
\endgroup
\noindent\textit{Key output:}
\begingroup\footnotesize
\begin{verbatim}
(ii)  IQ ~ nearc4                         n = 2061, R2 = 0.0059
             coef      se       t        p
Intercept  100.611   0.627
nearc4       2.596   0.745   3.483   0.0005
(iii) IQ ~ nearc4 + smsa66 + reg662..reg669          R2 = 0.0626
Intercept  104.773
nearc4       0.348   0.814   0.427   0.669
smsa66       1.089   0.809      (reg665 -8.475, reg667 -8.394)
F(9, 2050) for the geographic dummies = 13.80, p < 0.001;  corr(nearc4, smsa66) = 0.412
\end{verbatim}
\endgroup

\passo{(i) exogeneity}{Yes, it could. Families that value education, or parents with more schooling and
income, may choose to live near colleges; and colleges are in cities, where schools are better and
measured ability higher. Any such factor raises wages directly (it is in $u$) and is correlated with
$nearc4$, which would violate $\operatorname{Cov}(nearc4,u)=0$.}

\passo{(ii) unconditional}{The simple regression gives
\begin{equation}
\widehat{IQ}=100.611+2.596\,nearc4,\qquad t=\frac{(2.5962)}{(0.7455)}=3.4825,\label{q7:simple}
\end{equation}
so men who grew up near a four-year college have IQs about 2.6 points higher on average, a significant
difference ($p=0.0005$). Taken alone, this suggests $nearc4$ is correlated with ability.}

\passo{(iii) conditional}{With the geographic controls,
\begin{equation}
\widehat{IQ}=104.773+0.348\,nearc4+1.089\,smsa66+G_i'\hat\psi,\qquad t=\frac{(0.3479)}{(0.8144)}=0.4272,\label{q7:cond}
\end{equation}
with $p=0.669$: once geography is partialled out, IQ and $nearc4$ are unrelated. The geographic
dummies are jointly very significant ($F=13.80$). Reconciliation by the OVB formula of Question 2:
the short coefficient 2.596 equals the long coefficient 0.348 plus (effect of geography on IQ) $\times$
(relation of geography to $nearc4$). Men near colleges are more often urban ($\operatorname{corr}(nearc4,
smsa66)=0.412$) and live in regions with higher average IQ; that, not proximity, produced part (ii).}

\passo{(iv) implication}{$nearc4$ is a plausible instrument only \emph{conditional on} $smsa66$ and the
1966 regional dummies. They must be in the $\log(wage)$ equation (and hence in the first stage), as in
Table 15.1; dropping them would let geography-driven ability differences contaminate the IV estimate.}

\intuicao{Instrument validity is a conditional statement. A regression of a proxy for the omitted factor
(IQ) on the instrument, with and without the controls, is a cheap falsification test of that statement.}
\refbib{W \S15.1--15.2 and Example 15.4 (pp.~462--474); MHE \S4.1 (pp.~84--102); site example \texttt{example1-iv-card}.}
\end{sol}

% ---------------------------------------------------------------- Q8
\questao{Wooldridge, Computer Exercise 15.C5.}
\fonte{Wooldridge (2016), \textit{Introductory Econometrics: A Modern Approach}, 6th ed., Chapter 15,
Computer Exercise C5, printed p.~493, reproduced verbatim.}

\noindent\textbf{C5} Use the data in CARD for this exercise.
\begin{enumerate}[label=(\roman*)]
\item In Table 15.1, the difference between the IV and OLS estimates of the return to education is
economically important. Obtain the reduced form residuals, $\hat v_2$, from the reduced form regression
$educ$ on $nearc4$, $exper$, $exper^2$, $black$, $smsa$, $south$, $smsa66$, $reg662$, \dots,
$reg669$---see Table 15.1. Use these to test whether $educ$ is exogenous; that is, determine if the
difference between OLS and IV is statistically significant.
\item Estimate the equation by 2SLS, adding $nearc2$ as an instrument. Does the coefficient on $educ$ change much?
\item Test the single overidentifying restriction from part (ii).
\end{enumerate}
\resolespaco
\begin{sol}
\begin{notacao}
\item $\log(wage)_i$ ($lwage$): log hourly wage; $educ_i$: years of schooling; $x_i$: the exogenous
controls of Table 15.1 ($exper$, $exper^2$, $black$, $smsa$, $south$, $smsa66$, $reg662$--$reg669$).
\item $nearc4_i$, $nearc2_i$: grew up near a four-year / two-year college (instruments); $v_{2i}$:
reduced-form error; $\theta$: coefficient of $u_{1i}$ on $v_{2i}$.
\item $\hat u_i$: 2SLS residuals; $R^2_u$: $R^2$ of the overidentification regression; $n=3010$; $q$:
number of overidentifying restrictions.
\end{notacao}

\noindent\textit{Code:}
\begingroup\footnotesize
\begin{verbatim}
from linearmodels.iv import IV2SLS
X  = "exper+expersq+black+smsa+south+smsa66+" + regs
rf = smf.ols("educ ~ nearc4 + " + X, card).fit()                 # reduced form
card["v2hat"] = rf.resid
cf = smf.ols("lwage ~ educ + " + X + " + v2hat", card).fit()      # part (i)
iv2 = IV2SLS.from_formula("lwage ~ 1 + " + X + " + [educ ~ nearc4 + nearc2]",
                          card).fit(cov_type="unadjusted", debiased=True)   # part (ii)
card["u2sls"] = iv2.resids
ov = smf.ols("u2sls ~ nearc4 + nearc2 + " + X, card).fit()        # part (iii)
print(ov.nobs * ov.rsquared, iv2.sargan)
\end{verbatim}
\endgroup
\noindent\textit{Key output:}
\begingroup\footnotesize
\begin{verbatim}
Table 15.1 replicated: OLS educ 0.0747 (0.0035);  IV (nearc4) educ 0.1315 (0.0550)
(i)   reduced form: nearc4 0.320 (0.088), t = 3.64;  R2 = 0.477
      lwage ~ educ + X + v2hat:  educ 0.1315;  v2hat -0.0571 (0.0528), t = -1.081, p = 0.280
      heteroskedasticity-robust (HC1) t on v2hat = -1.101
(ii)  2SLS with nearc4, nearc2:  educ 0.1571 (0.0526)
      first stage: nearc4 0.321 (0.088), nearc2 0.123 (0.077)
(iii) R2 of u2sls on all exogenous variables = 0.00041467;  n*R2 = 1.248;  p = 0.264
      linearmodels Sargan statistic = 1.248, chi2(1)
\end{verbatim}
\endgroup

\passo{(i) control-function test}{Write the structural equation, the reduced form and the projection of
$u_1$ on $v_2$:
\begin{align}
lwage_i&=\beta_0+\beta_1educ_i+x_i'\beta+u_{1i},\label{q8:struct}\\
educ_i&=\pi_0+\pi_1nearc4_i+x_i'\pi+v_{2i},\label{q8:rf}\\
u_{1i}&=\theta\,v_{2i}+e_{1i}.\label{q8:proj}
\end{align}
Since $nearc4$ and $x_i$ are exogenous, $educ$ is correlated with $u_1$ only through $v_2$, so $educ$ is
exogenous if and only if $\theta=0$. Substituting \eqref{q8:proj} into \eqref{q8:struct} and replacing
$v_2$ by $\hat v_2$:
\begin{equation}
lwage_i=\beta_0+\beta_1educ_i+x_i'\beta+\theta\,\hat v_{2i}+\text{error}_i.\label{q8:cf}
\end{equation}
The $t$ statistic on $\hat v_2$ in \eqref{q8:cf} is
\begin{equation}
t=\frac{\hat\theta}{\widehat{se}(\hat\theta)}=\frac{(-0.05706)}{(0.05281)}=-1.0805,\label{q8:t}
\end{equation}
with $p=0.280$ (robust $t=-1.101$). We do not reject $H_0:\theta=0$: although the IV estimate (0.1315) is
almost twice OLS (0.0747), the difference is not statistically significant, because the IV estimate is
imprecise. (The coefficient on $educ$ in \eqref{q8:cf} equals the IV estimate, 0.1315, an algebraic identity.)}

\passo{(ii) adding nearc2}{2SLS with both instruments gives $\hat\beta_1=0.1571$ (0.0526). The change is
\begin{equation}
0.1571-0.1315=0.0256,\label{q8:change}
\end{equation}
about half a standard error: a moderate increase, with a slightly smaller standard error. $nearc2$ has a
smaller and less precise first-stage coefficient (0.123, se 0.077) than $nearc4$ (0.321, se 0.088), so it adds little.}

\passo{(iii) overidentification}{Regress the 2SLS residuals on all exogenous variables (both instruments
and $x_i$); under $H_0$ that all instruments are exogenous, $nR^2_u\sim\chi^2_q$ with
$q=2-1=1$ (instruments minus endogenous regressors). Substituting:
\begin{equation}
nR^2_u=(3010)(0.00041467)=1.24816,\label{q8:nr2}
\end{equation}
below the 5\% critical value $\chi^2_{1,0.95}=3.841$ ($p=0.264$). We do not reject the overidentifying
restriction: the two instruments give estimates (0.1315 alone, 0.1571 together) that are statistically
compatible. The test only checks that the instruments agree; it cannot detect a bias they share.}

\intuicao{The Hausman test asks whether OLS and IV differ by more than sampling noise; with a weak-ish
instrument the IV noise is large, so even a big point difference is not significant. The overidentification
test asks whether two instruments tell the same story.}
\refbib{W \S15.5 (pp.~481--484), testing for endogeneity and overidentifying restrictions; Table 15.1 (p.~474); MHE \S4.1 (pp.~84--102).}
\end{sol}

% =============================================================================
\bloco{3}{Difference-in-differences}{Wooldridge Ch.~13 \,\textbullet\, Questions 9--11}

% ---------------------------------------------------------------- Q9
\questao{Wooldridge, Problem 13.2.}
\fonte{Wooldridge (2016), \textit{Introductory Econometrics: A Modern Approach}, 6th ed., Chapter 13,
Problem 2, printed p.~425, reproduced verbatim. Equation (13.9), which the problem refers to, is
reproduced from p.~410 for reference.}

\noindent\textbf{2} Using the data in KIELMC, the following equations were estimated using the years 1978 and 1981:
\[
\begin{aligned}
\widehat{\log(price)}&=\underset{(.26)}{11.49}-\underset{(.058)}{.547}\,nearinc+\underset{(.080)}{.394}\,y81\cdot nearinc\\
n&=321,\ R^2=.220
\end{aligned}
\]
and
\[
\begin{aligned}
\widehat{\log(price)}&=\underset{(.27)}{11.18}+\underset{(.044)}{.563}\,y81-\underset{(.067)}{.403}\,y81\cdot nearinc\\
n&=321,\ R^2=.337.
\end{aligned}
\]
Compare the estimates on the interaction term $y81\cdot nearinc$ with those from equation (13.9). Why are
the estimates so different?

\smallskip
\noindent\textit{For reference, equation (13.9), p.~410:}
\[
\widehat{\log(price)}=\underset{(.31)}{11.29}+\underset{(.045)}{.457}\,y81-\underset{(.055)}{.340}\,nearinc-\underset{(.083)}{.063}\,y81\cdot nearinc,\qquad n=321,\ R^2=.409.\tag{13.9}
\]
\resolespaco
\begin{sol}
\begin{notacao}
\item $\log(price)_i$: log selling price of house $i$; $y81_i$: sold in 1981 (after the incinerator was
announced); $nearinc_i$: within three miles of the incinerator site.
\item $\mu_{gt}$: mean log price in cell $g\in\{\text{far},\text{near}\}$, $t\in\{78,81\}$;
$\beta_0,\delta_0,\beta_1,\delta_1$: the coefficients of (13.9) on the constant, $y81$, $nearinc$, $y81\cdot nearinc$.
\item $n_{gt}$: number of sales in cell $(g,t)$; in KIELMC $n_{\text{far},78}=123$, $n_{\text{far},81}=102$,
$n_{\text{near},78}=56$, $n_{\text{near},81}=40$.
\end{notacao}

\passo{cell means from (13.9)}{(13.9) has four parameters for four cells, so it fits the cell means exactly:
\begin{equation}
\mu_{\text{far},78}=\beta_0,\quad\mu_{\text{far},81}=\beta_0+\delta_0,\quad\mu_{\text{near},78}=\beta_0+\beta_1,\quad\mu_{\text{near},81}=\beta_0+\delta_0+\beta_1+\delta_1.\label{q9:cells}
\end{equation}}

\passo{first equation (no $y81$)}{Its regressors define three groups: far houses in both years (constant
only), near houses in 1978 (constant $+$ $nearinc$) and near houses in 1981 (all three). OLS on group
dummies fits group means, so the interaction coefficient is
\begin{equation}
\hat\theta_1=\mu_{\text{near},81}-\mu_{\text{near},78}.\label{q9:e1}
\end{equation}
Substituting \eqref{q9:cells} into \eqref{q9:e1}:
\begin{align}
\hat\theta_1&=(\beta_0+\delta_0+\beta_1+\delta_1)-(\beta_0+\beta_1)\label{q9:e1sub}\\
&=\delta_0+\delta_1=(.457)+(-.063)=.394,\label{q9:e1res}
\end{align}
exactly the reported .394. The intercept is the pooled far mean,
$11.29+.457\,(102/225)=11.4972\approx11.49$.}

\passo{second equation (no $nearinc$)}{Now the groups are all 1978 sales, far 1981 sales and near 1981
sales, so
\begin{equation}
\hat\theta_2=\mu_{\text{near},81}-\mu_{\text{far},81}.\label{q9:e2}
\end{equation}
Substituting \eqref{q9:cells} into \eqref{q9:e2}:
\begin{align}
\hat\theta_2&=(\beta_0+\delta_0+\beta_1+\delta_1)-(\beta_0+\delta_0)\label{q9:e2sub}\\
&=\beta_1+\delta_1=(-.340)+(-.063)=-.403,\label{q9:e2res}
\end{align}
again exactly the reported value. The intercept is the pooled 1978 mean,
$11.29-.340\,(56/179)=11.1836\approx11.18$.}

\passo{why they differ}{Each short equation imposes a false restriction. The first sets $\delta_0=0$ (no
general change in prices from 1978 to 1981), so its interaction is the change in prices of near houses,
which mixes the aggregate appreciation ($+.457$) with the incinerator effect. The second sets $\beta_1=0$
(near and far houses equally priced before the incinerator), so its interaction is the 1981 near--far gap,
which mixes the pre-existing location discount ($-.340$) with the effect. Only (13.9), with both controls,
isolates the difference-in-differences $\hat\delta_1=-.063$. (Fitting the three equations to KIELMC
reproduces all coefficients; the intercept standard errors are .026 and .027 in the data, so the printed
(.26) and (.27) appear to have lost a zero.)}

\intuicao{A DiD regression needs both a group dummy and a period dummy; dropping either one lets the
interaction soak up the omitted level difference, a textbook case of omitted-variable bias with dummies.}
\refbib{W Example 13.3 (pp.~407--410) and Table 13.3 (p.~411); MHE \S5.2 (pp.~169--171).}
\end{sol}

% ---------------------------------------------------------------- Q10
\questao{Wooldridge, Problem 13.6.}
\fonte{Wooldridge (2016), \textit{Introductory Econometrics: A Modern Approach}, 6th ed., Chapter 13,
Problem 6, printed p.~426, reproduced verbatim.}

\noindent\textbf{6} In 1985, neither Florida nor Georgia had laws banning open alcohol containers in vehicle
passenger compartments. By 1990, Florida had passed such a law, but Georgia had not.
\begin{enumerate}[label=(\roman*)]
\item Suppose you can collect random samples of the driving-age population in both states, for 1985 and
1990. Let $arrest$ be a binary variable equal to unity if a person was arrested for drunk driving during
the year. Without controlling for any other factors, write down a linear probability model that allows you
to test whether the open container law reduced the probability of being arrested for drunk driving. Which
coefficient in your model measures the effect of the law?
\item Why might you want to control for other factors in the model? What might some of these factors be?
\item Now, suppose that you can only collect data for 1985 and 1990 at the county level for the two
states. The dependent variable would be the fraction of licensed drivers arrested for drunk driving during
the year. How does this data structure differ from the individual-level data described in part (i)? What
econometric method would you use?
\end{enumerate}
\resolespaco
\begin{sol}
\begin{notacao}
\item $arrest_i\in\{0,1\}$: person $i$ arrested for drunk driving in the year; $y90_i$: observed in 1990;
$FL_i$: lives in Florida (treated state); Georgia is the control.
\item $\delta_1$: effect of the open-container law on the arrest probability; $x$: other factors;
$arrrate_{ct}$: fraction of licensed drivers arrested in county $c$, year $t$; $a_c$: county effect.
\end{notacao}

\passo{(i) the model}{The linear probability model is
\begin{equation}
arrest_i=\beta_0+\delta_0\,y90_i+\beta_1\,FL_i+\delta_1\,y90_i\cdot FL_i+u_i.\label{q10:lpm}
\end{equation}
Taking expectations of \eqref{q10:lpm} in the four state--year cells and differencing as in Question 4:
\begin{equation}
\delta_1=\big[P(FL,90)-P(FL,85)\big]-\big[P(GA,90)-P(GA,85)\big],\label{q10:did}
\end{equation}
where $P(s,t)$ is the arrest probability in state $s$, year $t$. $\delta_1$ measures the effect of the law;
test $H_0:\delta_1=0$ against $H_1:\delta_1<0$ with the $t$ statistic on the interaction
(heteroskedasticity-robust, since an LPM is heteroskedastic).}

\passo{(ii) controls}{The samples are independent cross sections, so their composition can change
differently in the two states between 1985 and 1990; then \eqref{q10:did} would mix the law's effect with
composition effects. Candidates: age, sex, income, education, urban residence, and state-level factors that
may have changed differently, such as police enforcement and sobriety checkpoints, beer taxes, the legal
blood-alcohol limit or the drinking age. Controls that do not change differently still help by lowering the
residual variance and the standard error of $\hat\delta_1$.}

\passo{(iii) county data}{Now the unit is the county, observed in both years: a two-period panel, and the
outcome is a proportion. Write
\begin{equation}
arrrate_{ct}=\beta_0+\delta_0\,y90_t+\delta_1\,y90_t\cdot FL_c+x_{ct}'\gamma+a_c+u_{ct},\label{q10:panel}
\end{equation}
where $a_c$ is a time-constant county effect. First-differencing \eqref{q10:panel} between 1990 and 1985:
\begin{equation}
\Delta arrrate_c=\delta_0+\delta_1\,FL_c+\Delta x_c'\gamma+\Delta u_c,\label{q10:fd}
\end{equation}
which removes $a_c$ and is estimated by OLS on one observation per county (first-differencing, equal to
fixed effects with two periods). Because the outcome is a fraction, weight by the number of licensed
drivers, and since the law varies only at the state level with just two states, standard errors should
allow for state-level correlation and be read with caution.}

\intuicao{The individual design compares the change in Florida with the change in Georgia across
different people; the county design follows the same counties, so differencing removes their permanent
differences.}
\refbib{W \S13.2 (pp.~406--412), policy analysis with pooled cross sections; \S13.3--13.4, two-period panels and first differencing.}
\end{sol}

% ---------------------------------------------------------------- Q11
\questao{Wooldridge, Computer Exercise 13.C3.}
\fonte{Wooldridge (2016), \textit{Introductory Econometrics: A Modern Approach}, 6th ed., Chapter 13,
Computer Exercise C3, printed p.~427, reproduced verbatim.}

\noindent\textbf{C3} Use the data in KIELMC for this exercise.
\begin{enumerate}[label=(\roman*)]
\item The variable $dist$ is the distance from each home to the incinerator site, in feet. Consider the model
\[\log(price)=\beta_0+\delta_0y81+\beta_1\log(dist)+\delta_1y81\cdot\log(dist)+u.\]
If building the incinerator reduces the value of homes closer to the site, what is the sign of $\delta_1$?
What does it mean if $\beta_1>0$?
\item Estimate the model from part (i) and report the results in the usual form. Interpret the coefficient
on $y81\cdot\log(dist)$. What do you conclude?
\item Add $age$, $age^2$, $rooms$, $baths$, $\log(intst)$, $\log(land)$, and $\log(area)$ to the equation.
Now, what do you conclude about the effect of the incinerator on housing values?
\item Why is the coefficient on $\log(dist)$ positive and statistically significant in part (ii) but not in
part (iii)? What does this say about the controls used in part (iii)?
\end{enumerate}
\resolespaco
\begin{sol}
\begin{notacao}
\item $\log(price)_i$: log selling price; $\log(dist)_i$: log distance to the incinerator site (feet);
$y81_i$: sold in 1981.
\item $\beta_1$: price--distance elasticity in 1978; $\beta_1+\delta_1$: the elasticity in 1981;
$intst$: distance to the interstate; $land$, $area$: lot and house area (square feet).
\end{notacao}

\noindent\textit{Code} (KIELMC from the \texttt{wooldridge} package, saved as \texttt{Listas/data/kielmc.csv}):
\begingroup\footnotesize
\begin{verbatim}
k  = pd.read_csv("Listas/data/kielmc.csv")
m2 = smf.ols("lprice ~ y81 + ldist + y81ldist", k).fit()                          # part (ii)
m3 = smf.ols("lprice ~ y81 + ldist + y81ldist + age + agesq + rooms + baths"
             " + lintst + lland + larea", k).fit()                                # part (iii)
\end{verbatim}
\endgroup
\noindent\textit{Key output:}
\begingroup\footnotesize
\begin{verbatim}
(ii)  lprice ~ y81 + ldist + y81ldist          n = 321, R2 = 0.396
             coef      se       t        p
Intercept   8.058    0.508
y81        -0.011    0.805
ldist       0.317    0.052    6.15
y81ldist    0.048    0.082    0.589    0.556
(iii) + age, agesq, rooms, baths, lintst, lland, larea          R2 = 0.787
y81        -0.2254   0.4947
ldist       0.0009   0.0446   0.021
y81ldist    0.0625   0.0503   1.242    0.215
lintst     -0.0600;  larea 0.3507;  corr(ldist, lintst) = 0.781
\end{verbatim}
\endgroup

\passo{(i) signs}{From the model, the price--distance elasticity is
\begin{equation}
\frac{\partial\log(price)}{\partial\log(dist)}=\beta_1+\delta_1\,y81,\label{q11:elas}
\end{equation}
equal to $\beta_1$ in 1978 and $\beta_1+\delta_1$ in 1981. If the incinerator lowers the value of nearby
homes, being farther away becomes worth more after it is built, so the elasticity rises: $\delta_1>0$.
$\beta_1>0$ means that already in 1978, before the incinerator was rumored, homes farther from the future
site sold for more: the site was in a less valuable area (as Example 13.3 found with $nearinc$).}

\passo{(ii) estimates}{The results in the usual form:
\begin{equation}
\widehat{\log(price)}=\underset{(0.508)}{8.058}-\underset{(0.805)}{0.011}\,y81+\underset{(0.052)}{0.317}\log(dist)+\underset{(0.082)}{0.048}\,y81\cdot\log(dist),\quad n=321,\ R^2=0.396.\label{q11:ii}
\end{equation}
By \eqref{q11:elas}, after the incinerator a 1\% greater distance adds about 0.048\% more to the price than
before. The $t$ statistic is $(0.0482)/(0.0818)=0.589$ ($p=0.556$): the sign is as predicted but there is no
statistical evidence of an incinerator effect.}

\passo{(iii) with controls}{With the house characteristics, $\hat\delta_1=0.0625$ (0.0503), $t=1.242$,
$p=0.215$. The estimate is larger but still not significant at conventional levels, so this specification,
which uses distance rather than the \emph{nearinc} dummy, gives no firm evidence that the incinerator lowered
housing values (contrast the $-.132$, $t\approx-2.53$, that Wooldridge reports with $nearinc$ and the same controls).}

\passo{(iv) the role of the controls}{In (ii) $\log(dist)$ has coefficient 0.317 with $t=6.15$; in (iii)
0.0009 with $t=0.021$. Homes far from the site differ systematically: they are larger, on larger lots, newer,
and farther from the interstate ($\operatorname{corr}(\log dist,\log intst)=0.781$). In (ii) $\log(dist)$ picks
up these characteristics (omitted-variable bias); in (iii) they are controlled for and the pure location
effect of distance vanishes. So the controls capture what made distant homes more expensive before the
incinerator, which is exactly why they belong in the DiD regression.}

\intuicao{A continuous treatment intensity (distance) replaces the near/far dummy; the DiD logic is the
same, but now the identifying variation is the change in the price--distance gradient.}
\refbib{W Example 13.3 (pp.~407--410); Kiel and McClain (1995); MHE \S5.2 (pp.~169--181), regression DD with covariates.}
\end{sol}

% =============================================================================
\bloco{4}{Bayesian treatment models}{Greenberg Ch.~11 \,\textbullet\, Questions 12--15}

\noindent\textit{Context for Questions 12--15} (Greenberg 2007, \S11.1, printed pp.~169--172; the model and
the results the exercises refer to, reproduced verbatim, selected passages). The response is related to the
covariates and the treatment assignment through
\[
y_{0i}=x_i'\beta_0+u_{0i},\qquad y_{1i}=x_i'\beta_1+u_{1i},\tag{11.1}
\]
``where $x_i$ is a $K_1$ vector of covariates, a 0 subscript indicates assignment to the control group, and
1 indicates assignment to the treatment group.'' Assignment is modeled by
\[
s_i^*=x_i'\gamma_1+z_i'\gamma_2+v_i,\qquad s_i=\begin{cases}1,&\text{if }s_i^*>0,\\0,&\text{if }s_i^*\leq0,\end{cases}\tag{11.2}
\]
\[
y_i=(1-s_i)y_{0i}+s_iy_{1i},\tag{11.3}
\]
``where $z_i$ is a $K_2\times1$ vector of instrumental variables.''
\[
\operatorname{Cov}(u_{0i},v_i)\equiv\Sigma_0=\begin{pmatrix}\sigma_{00}&\sigma_{0v}\\\sigma_{0v}&1\end{pmatrix},\qquad
\operatorname{Cov}(u_{1i},v_i)\equiv\Sigma_1=\begin{pmatrix}\sigma_{11}&\sigma_{1v}\\\sigma_{1v}&1\end{pmatrix},
\]
``where $\sigma_{vv}=1$ because of the probit specification.'' ``Let $\beta=(\beta_0',\beta_1',\gamma_1',\gamma_2')'$,
$\gamma=(\gamma_1',\gamma_2')'$,
$X_{ji}=\begin{pmatrix}(1-j)x_i'&jx_i'&0&0\\0&0&x_i'&z_i'\end{pmatrix}$, [\dots] $N_0=\{i:s_i=0\}$,
$N_1=\{i:s_i=1\}$, and $n_0$ and $n_1$, respectively, are the number of observations in $N_0$ and $N_1$. We
assume a Gaussian prior distribution for $\beta\sim N_p(b_0,B_0)$, where $p=3K_1+K_2$. [\dots] we
reparameterize the variances in $\Sigma_j$ as $\omega_{jj}=\sigma_{jj}-\sigma_{jv}^2$ [\dots]. We define
$\psi_j=(\omega_{jj},\sigma_{jv})$ and specify the prior distributions
$\pi(\psi_j)=\pi(\omega_{jj})\pi(\sigma_{jv}\mid\omega_{jj})=\mathrm{IG}(\nu_{j0}/2,d_{j0}/2)N(m_{j0},\omega_{jj}M_{j0})$.''
\[
y_{ji}\mid s_i^*,\beta_j,\gamma,\psi_j\sim N(x_i'\beta_j+\sigma_{jv}e_i,\omega_{jj}),\tag{11.4}
\]
``where $e_i=s_i^*-x_i'\gamma_1-z_i'\gamma_2$. Let $Y_j$, $X_j$, $Z_j$, $S_j^*$, and $E_j$, respectively, contain
the values of $y_{ji}$, $x_i$, $z_i$, $s_i^*$, and $e_i$ for $i\in N_j$. Then''
\[
Y_j\mid S_j^*,\beta_j,\gamma,\psi_j\sim N_{n_j}\big(X_j\beta_j+\sigma_{jv}E_j,\ \omega_{jj}I_{n_j}\big).\tag{11.5}
\]
\[
E(Y_j\mid S_j^*,\beta_j,\gamma,\omega_{jj})=X_j\beta_j+m_{j0}E_j,\tag{11.6}
\]
\[
\operatorname{Var}(Y_j\mid S_j^*,\gamma,\omega_{jj})=\omega_{jj}\big(I_{n_j}+E_jM_{j0}E_j'\big).\tag{11.7}
\]
\[
f(y_{ji},s_i=j\mid x_i,z_i,\beta,\psi_j)=N(y_{ji}\mid x_i'\beta_j,\sigma_{jj})\,\Phi\!\left((2j-1)\frac{\mu_{ji}}{\phi_j}\right).\tag{11.8}
\]
\textbf{Algorithm 11.1: Gibbs algorithm for treatment model (Chib, in press-a)}
\begin{description}[leftmargin=2.2em,font=\normalfont\bfseries]
\item[11.1] Sample $\psi_j$ conditioned on $(y,s^*,X,Z,\beta)$\\
(a) Drawing $\omega_{jj}$ marginalized over $\sigma_{jv}$ from
$\mathrm{IG}\left(\frac{\nu_{j0}+n_j}{2},\frac{d_{j0}+d_j}{2}\right)$, where
$d_j=(Y_j-X_j\beta_j-m_{j0}E_j)'\big(I_{n_j}+E_jM_{j0}E_j'\big)^{-1}(Y_j-X_j\beta_j-m_{j0}E_j)$,\\
(b) Drawing $\sigma_{jv}$ conditioned on $\omega_{jj}$ from $N(c_j,\omega_{jj}C_j)$, where
$c_j=C_j\big[M_{j0}^{-1}m_{j0}+E_j'(Y_j-X_j\beta_j)\big]$, $C_j=\big(M_{j0}^{-1}+E_j'E_j\big)^{-1}$.
\item[11.2] Sample $\beta$ conditioned on $(y,s^*,\psi_0,\psi_1)$ from $N_p(\hat\beta,B)$, where
$\hat\beta=B\big(B_0^{-1}b_0+\sum_{i\in N_0}X_{0i}'\Sigma_0^{-1}y_i^*+\sum_{i\in N_1}X_{1i}'\Sigma_1^{-1}y_i^*\big)$,
$B=\big(B_0^{-1}+\sum_{i\in N_0}X_{0i}'\Sigma_0^{-1}X_{0i}+\sum_{i\in N_1}X_{1i}'\Sigma_1^{-1}X_{1i}\big)^{-1}$,
and $y_i^*=(y_i,s_i^*)'$.
\item[11.3] Sample $s^*$ conditioned on $(y,X,Z,\beta,\psi_0,\psi_1)$ by drawing $s_i^*$ from
$\mathrm{TN}_{(-\infty,0)}(\mu_{0i},\phi_0^2)$ if $s_i=0$, $\mathrm{TN}_{(0,\infty)}(\mu_{1i},\phi_1^2)$ if $s_i=1$, where
$\mu_{ji}=x_i'\gamma_1+z_i'\gamma_2+\sigma_{jv}\sigma_{jj}^{-1}(y_{ji}-x_i'\beta_j)$, $\phi_j^2=1-\sigma_{jv}^2\sigma_{jj}^{-1}$.
\end{description}

\noindent\textit{Common notation for the solutions.} $j\in\{0,1\}$: untreated/treated; $\Phi$: standard
normal c.d.f.; $\mathrm{IG}(\alpha,\beta)$ has density proportional to $\omega^{-(\alpha+1)}e^{-\beta/\omega}$
($\omega>0$); $\mathrm{TN}_{(a,b)}$: normal truncated to $(a,b)$. The bivariate-normal conditioning rule
(Greenberg, Appendix A, eq.~A.9) is used throughout: if $(a,b)'$ is bivariate normal,
\[
a\mid b\sim N\!\left(\mu_a+\frac{\sigma_{ab}}{\sigma_{bb}}(b-\mu_b),\ \sigma_{aa}-\frac{\sigma_{ab}^2}{\sigma_{bb}}\right).\tag{A.9}
\]

% ---------------------------------------------------------------- Q12
\questao{Greenberg, Exercise 11.1.}
\fonte{Greenberg (2007), \textit{Introduction to Bayesian Econometrics}, 1st ed., Chapter 11, Exercise
11.1, printed p.~180, reproduced verbatim.}

\noindent\textbf{11.1} Verify (11.4).
\resolespaco
\begin{sol}
\begin{notacao}
\item $u_{ji},v_i$: errors of the outcome and assignment equations, jointly $N_2(0,\Sigma_j)$;
$\sigma_{jj}>0$, $\sigma_{jv}$ real, $\omega_{jj}=\sigma_{jj}-\sigma_{jv}^2>0$.
\item $e_i=s_i^*-x_i'\gamma_1-z_i'\gamma_2$: the value of $v_i$ implied by $s_i^*$ and $\gamma$.
\end{notacao}

\passo{the ingredients}{From (11.1), (11.2) and the covariance matrix $\Sigma_j$:
\begin{align}
y_{ji}&=x_i'\beta_j+u_{ji},\label{q12:y}\\
s_i^*&=x_i'\gamma_1+z_i'\gamma_2+v_i,\label{q12:s}\\
\begin{pmatrix}u_{ji}\\v_i\end{pmatrix}&\sim N_2\!\left(\begin{pmatrix}0\\0\end{pmatrix},\begin{pmatrix}\sigma_{jj}&\sigma_{jv}\\\sigma_{jv}&1\end{pmatrix}\right),\label{q12:joint}\\
\sigma_{jj}&=\omega_{jj}+\sigma_{jv}^2.\label{q12:omega}
\end{align}}

\passo{conditioning on $s^*$ fixes $v$}{Given $s_i^*$ and $\gamma$, solving \eqref{q12:s} for $v_i$:
\begin{equation}
v_i=s_i^*-x_i'\gamma_1-z_i'\gamma_2=e_i.\label{q12:e}
\end{equation}
So conditioning on $(s_i^*,\gamma)$ is conditioning on $v_i=e_i$.}

\passo{conditional distribution of $u$}{Apply (A.9) with $a=u_{ji}$, $b=v_i$, and from \eqref{q12:joint}
$\mu_a=\mu_b=0$, $\sigma_{aa}=\sigma_{jj}$, $\sigma_{ab}=\sigma_{jv}$, $\sigma_{bb}=1$. The mean is
\begin{align}
E(u_{ji}\mid v_i)&=(0)+\frac{(\sigma_{jv})}{(1)}\big(v_i-(0)\big)\label{q12:m1}\\
&=\sigma_{jv}e_i,\label{q12:m2}
\end{align}
using \eqref{q12:e}. The variance is
\begin{align}
\operatorname{Var}(u_{ji}\mid v_i)&=(\sigma_{jj})-\frac{(\sigma_{jv})^2}{(1)}\label{q12:v1}\\
&=\omega_{jj},\label{q12:v2}
\end{align}
by \eqref{q12:omega}.}

\passo{back to $y$}{Given $(s_i^*,\beta_j,\gamma,\psi_j)$, $x_i'\beta_j$ in \eqref{q12:y} is a constant, so
$y_{ji}$ is $u_{ji}$ shifted by $x_i'\beta_j$. Substituting \eqref{q12:m2} and \eqref{q12:v2}:
\begin{equation}
y_{ji}\mid s_i^*,\beta_j,\gamma,\psi_j\sim N\big(x_i'\beta_j+(\sigma_{jv}e_i),\ (\omega_{jj})\big),\label{q12:res}
\end{equation}
which is (11.4). The distribution depends on $\psi_j=(\omega_{jj},\sigma_{jv})$ only, which is why the
reparameterization was chosen. Independence across $i$ then stacks \eqref{q12:res} into (11.5).}

\intuicao{Once the latent assignment index is known, the assignment error $v_i$ is known, and the part of
the outcome error correlated with it, $\sigma_{jv}e_i$, becomes a regressor: this is a Bayesian control
function, the analogue of adding $\hat v_2$ in Question 8.}
\refbib{Greenberg \S11.1 (pp.~168--172) and Appendix A (eq.~A.9); compare the control-function test of W \S15.5 (Question 8).}
\end{sol}

% ---------------------------------------------------------------- Q13
\questao{Greenberg, Exercise 11.2.}
\fonte{Greenberg (2007), \textit{Introduction to Bayesian Econometrics}, 1st ed., Chapter 11, Exercise
11.2, printed p.~180, reproduced verbatim.}

\noindent\textbf{11.2} Verify (11.6) and (11.7).
\resolespaco
\begin{sol}
\begin{notacao}
\item $Y_j,X_j,E_j$: stacked $y_{ji}$, $x_i'$, $e_i$ for $i\in N_j$ ($n_j\times1$, $n_j\times K_1$, $n_j\times1$).
\item $m_{j0}$, $M_{j0}>0$: prior mean and scale of $\sigma_{jv}$ given $\omega_{jj}$; $\eta_j$: prior
deviation of $\sigma_{jv}$ from $m_{j0}$; $\varepsilon_j$: the $n_j$-vector of conditional errors in (11.5).
\end{notacao}

\passo{model and prior}{(11.5) and the prior on $\sigma_{jv}$ are
\begin{align}
Y_j&=X_j\beta_j+\sigma_{jv}E_j+\varepsilon_j,\qquad\varepsilon_j\sim N_{n_j}(0,\omega_{jj}I_{n_j}),\label{q13:model}\\
\sigma_{jv}&=m_{j0}+\eta_j,\qquad\eta_j\sim N(0,\omega_{jj}M_{j0}),\label{q13:prior}
\end{align}
with $\eta_j$ independent of $\varepsilon_j$ given $\omega_{jj}$ (prior information versus sampling noise).}

\passo{marginalize $\sigma_{jv}$}{Substituting \eqref{q13:prior} into \eqref{q13:model}:
\begin{equation}
Y_j=X_j\beta_j+(m_{j0}+\eta_j)E_j+\varepsilon_j.\label{q13:sub}
\end{equation}
Separating the fixed and random parts of \eqref{q13:sub}:
\begin{equation}
Y_j=X_j\beta_j+m_{j0}E_j+\big(\eta_jE_j+\varepsilon_j\big).\label{q13:split}
\end{equation}
Given $(S_j^*,\gamma)$, $E_j$ is fixed, so $\eta_jE_j+\varepsilon_j$ is a linear combination of independent
normals and is normal.}

\passo{the mean}{Taking expectations of the bracket in \eqref{q13:split}:
\begin{equation}
E\big(\eta_jE_j+\varepsilon_j\big)=E_j\,E(\eta_j)+E(\varepsilon_j)=E_j(0)+(0)=0,\label{q13:mean}
\end{equation}
so $E(Y_j\mid S_j^*,\beta_j,\gamma,\omega_{jj})=X_j\beta_j+m_{j0}E_j$, which is (11.6).}

\passo{the variance}{By independence of $\eta_j$ and $\varepsilon_j$,
\begin{align}
\operatorname{Var}\big(\eta_jE_j+\varepsilon_j\big)&=E_j\operatorname{Var}(\eta_j)E_j'+\operatorname{Var}(\varepsilon_j)\label{q13:v1}\\
&=E_j(\omega_{jj}M_{j0})E_j'+(\omega_{jj}I_{n_j})\label{q13:v2}\\
&=\omega_{jj}\big(I_{n_j}+E_jM_{j0}E_j'\big),\label{q13:v3}
\end{align}
substituting the variances from \eqref{q13:model}--\eqref{q13:prior} in \eqref{q13:v2} and factoring
$\omega_{jj}$ in \eqref{q13:v3}. This is (11.7). (The variance does not involve $\beta_j$, which is why the
book omits it from the conditioning set in (11.7).)}

\intuicao{Integrating out $\sigma_{jv}$ turns its prior uncertainty into extra, correlated noise along the
direction $E_j$: observations with large $|e_i|$ get more variance because the control-function term is
uncertain.}
\refbib{Greenberg \S11.1 (pp.~170--171); conjugate normal--inverse-gamma regression, Greenberg Ch.~4.}
\end{sol}

% ---------------------------------------------------------------- Q14
\questao{Greenberg, Exercise 11.3.}
\fonte{Greenberg (2007), \textit{Introduction to Bayesian Econometrics}, 1st ed., Chapter 11, Exercise
11.3, printed p.~181, reproduced verbatim.}

\noindent\textbf{11.3} Verify (11.8).
\resolespaco
\begin{sol}
\begin{notacao}
\item $f(y_{ji},s_i=j\mid\cdot)$: joint density of the observed outcome and probability of the observed
assignment for unit $i$ in group $j$; $N(y\mid m,s^2)$: normal density at $y$.
\item $\mu_{ji}$, $\phi_j^2$: conditional mean and variance of $s_i^*$ given $y_{ji}$ (defined in step 11.3);
$\Phi(-a)=1-\Phi(a)$.
\end{notacao}

\passo{factor the joint}{Write the joint as marginal times conditional:
\begin{equation}
f(y_{ji},s_i=j\mid\cdot)=f(y_{ji}\mid\cdot)\,P(s_i=j\mid y_{ji},\cdot).\label{q14:fac}
\end{equation}}

\passo{the marginal of $y$}{From \eqref{q12:y} and \eqref{q12:joint} in Question 12,
$u_{ji}\sim N(0,\sigma_{jj})$, so
\begin{equation}
f(y_{ji}\mid\cdot)=N(y_{ji}\mid x_i'\beta_j,\sigma_{jj}).\label{q14:marg}
\end{equation}}

\passo{$s^*$ given $y$}{$(y_{ji},s_i^*)$ is bivariate normal with means $(x_i'\beta_j,\ x_i'\gamma_1+z_i'\gamma_2)$,
variances $(\sigma_{jj},1)$ and covariance $\operatorname{Cov}(u_{ji},v_i)=\sigma_{jv}$. Apply (A.9) with
$a=s_i^*$, $b=y_{ji}$:
\begin{align}
E(s_i^*\mid y_{ji})&=(x_i'\gamma_1+z_i'\gamma_2)+\frac{(\sigma_{jv})}{(\sigma_{jj})}\big(y_{ji}-x_i'\beta_j\big)=\mu_{ji},\label{q14:mu}\\
\operatorname{Var}(s_i^*\mid y_{ji})&=(1)-\frac{(\sigma_{jv})^2}{(\sigma_{jj})}=\phi_j^2,\label{q14:phi}
\end{align}
the definitions in step 11.3 of Algorithm 11.1.}

\passo{the probabilities}{For $j=1$, standardizing with \eqref{q14:mu}--\eqref{q14:phi}:
\begin{align}
P(s_i=1\mid y_{1i})&=P(s_i^*>0\mid y_{1i})=1-\Phi\!\left(\frac{0-\mu_{1i}}{\phi_1}\right)\label{q14:p1a}\\
&=\Phi\!\left(\frac{\mu_{1i}}{\phi_1}\right),\label{q14:p1b}
\end{align}
by $1-\Phi(-a)=\Phi(a)$. For $j=0$:
\begin{equation}
P(s_i=0\mid y_{0i})=P(s_i^*\leq0\mid y_{0i})=\Phi\!\left(\frac{0-\mu_{0i}}{\phi_0}\right)=\Phi\!\left(-\frac{\mu_{0i}}{\phi_0}\right).\label{q14:p0}
\end{equation}
Since $2j-1=1$ for $j=1$ and $2j-1=-1$ for $j=0$, \eqref{q14:p1b} and \eqref{q14:p0} are both
\begin{equation}
P(s_i=j\mid y_{ji})=\Phi\!\left((2j-1)\frac{\mu_{ji}}{\phi_j}\right).\label{q14:both}
\end{equation}}

\passo{assemble}{Substituting \eqref{q14:marg} and \eqref{q14:both} into \eqref{q14:fac}:
\begin{equation}
f(y_{ji},s_i=j\mid x_i,z_i,\beta,\psi_j)=\big(N(y_{ji}\mid x_i'\beta_j,\sigma_{jj})\big)\Big(\Phi\!\left((2j-1)\tfrac{\mu_{ji}}{\phi_j}\right)\Big),\label{q14:res}
\end{equation}
which is (11.8); for $j=0$ it is the book's $N(y_{0i}\mid x_i'\beta_0,\sigma_{00})\Phi(-\mu_{0i}/\phi_0)$.}

\intuicao{The likelihood contribution is ``probability of the outcome'' times ``probit probability of the
observed assignment, updated by what the outcome reveals about the assignment error''; with $\sigma_{jv}=0$
it splits into a regression and an ordinary probit.}
\refbib{Greenberg \S11.1 (p.~172), marginal likelihood by Chib's method; \S8.2.2 (binary probit with latent data).}
\end{sol}

% ---------------------------------------------------------------- Q15
\questao{Greenberg, Exercise 11.4.}
\fonte{Greenberg (2007), \textit{Introduction to Bayesian Econometrics}, 1st ed., Chapter 11, Exercise
11.4, printed p.~181, reproduced verbatim.}

\noindent\textbf{11.4} Verify the conditional posterior distributions used in Algorithm 11.1.
\resolespaco
\begin{sol}
\begin{notacao}
\item $r_j=Y_j-X_j\beta_j-m_{j0}E_j$: residual from the marginal mean (11.6); $A_j=I_{n_j}+E_jM_{j0}E_j'$;
$w_j=Y_j-X_j\beta_j$.
\item $\epsilon_i=y_i^*-X_{ji}\beta$: the $2\times1$ error of unit $i\in N_j$, $\epsilon_i\sim N_2(0,\Sigma_j)$;
$\propto$: proportional to, as a function of the parameter being sampled.
\end{notacao}

\passo{(a) $\omega_{jj}$ marginalized over $\sigma_{jv}$}{By Bayes, using the marginal density of $Y_j$ from
Question 13 and the IG prior:
\begin{equation}
\pi(\omega_{jj}\mid\cdot)\propto N_{n_j}\big(Y_j\mid X_j\beta_j+m_{j0}E_j,\ \omega_{jj}A_j\big)\ \pi(\omega_{jj}).\label{q15:a0}
\end{equation}
The normal density is proportional to $|\omega_{jj}A_j|^{-1/2}\exp\{-r_j'(\omega_{jj}A_j)^{-1}r_j/2\}$, and
$|\omega_{jj}A_j|=\omega_{jj}^{n_j}|A_j|$, so
\begin{equation}
N_{n_j}(\cdot)\propto\omega_{jj}^{-n_j/2}\exp\left\{-\frac{d_j}{2\omega_{jj}}\right\},\qquad d_j=r_j'A_j^{-1}r_j.\label{q15:a1}
\end{equation}
The prior $\mathrm{IG}(\nu_{j0}/2,d_{j0}/2)$ is
\begin{equation}
\pi(\omega_{jj})\propto\omega_{jj}^{-(\nu_{j0}/2+1)}\exp\left\{-\frac{d_{j0}}{2\omega_{jj}}\right\}.\label{q15:a2}
\end{equation}
Substituting \eqref{q15:a1} and \eqref{q15:a2} into \eqref{q15:a0}:
\begin{align}
\pi(\omega_{jj}\mid\cdot)&\propto\left(\omega_{jj}^{-n_j/2}e^{-d_j/(2\omega_{jj})}\right)\left(\omega_{jj}^{-(\nu_{j0}/2+1)}e^{-d_{j0}/(2\omega_{jj})}\right)\label{q15:a3}\\
&=\omega_{jj}^{-\left(\frac{\nu_{j0}+n_j}{2}+1\right)}\exp\left\{-\frac{d_{j0}+d_j}{2\,\omega_{jj}}\right\},\label{q15:a4}
\end{align}
the kernel of $\mathrm{IG}\big((\nu_{j0}+n_j)/2,(d_{j0}+d_j)/2\big)$, with $d_j$ as in step 11.1(a).}

\passo{(b) $\sigma_{jv}$ given $\omega_{jj}$}{The likelihood (11.5) and the prior give
\begin{equation}
\pi(\sigma_{jv}\mid\omega_{jj},\cdot)\propto\exp\left\{-\frac{(w_j-\sigma_{jv}E_j)'(w_j-\sigma_{jv}E_j)}{2\omega_{jj}}\right\}\exp\left\{-\frac{(\sigma_{jv}-m_{j0})^2}{2\omega_{jj}M_{j0}}\right\}.\label{q15:b0}
\end{equation}
Multiplying the exponent of \eqref{q15:b0} by $-2\omega_{jj}$ and expanding:
\begin{equation}
\sigma_{jv}^2E_j'E_j-2\sigma_{jv}E_j'w_j+w_j'w_j+M_{j0}^{-1}\big(\sigma_{jv}^2-2\sigma_{jv}m_{j0}+m_{j0}^2\big).\label{q15:b1}
\end{equation}
Collecting powers of $\sigma_{jv}$ in \eqref{q15:b1}:
\begin{equation}
\sigma_{jv}^2\big(M_{j0}^{-1}+E_j'E_j\big)-2\sigma_{jv}\big(M_{j0}^{-1}m_{j0}+E_j'w_j\big)+\text{const}.\label{q15:b2}
\end{equation}
With $C_j=(M_{j0}^{-1}+E_j'E_j)^{-1}$ and $c_j=C_j(M_{j0}^{-1}m_{j0}+E_j'w_j)$, \eqref{q15:b2} is
\begin{equation}
C_j^{-1}\big(\sigma_{jv}-c_j\big)^2+\text{const},\label{q15:b3}
\end{equation}
so $\pi(\sigma_{jv}\mid\omega_{jj},\cdot)\propto\exp\{-(\sigma_{jv}-c_j)^2/(2\omega_{jj}C_j)\}$, i.e.\
$N(c_j,\omega_{jj}C_j)$ with $w_j=Y_j-X_j\beta_j$, as in step 11.1(b).}

\passo{(c) $\beta$}{For $i\in N_j$ stack $y_i^*=(y_i,s_i^*)'$. By the definition of $X_{ji}$,
\begin{equation}
X_{ji}\beta=\begin{pmatrix}(1-j)x_i'\beta_0+jx_i'\beta_1\\x_i'\gamma_1+z_i'\gamma_2\end{pmatrix},\qquad y_i^*=X_{ji}\beta+\epsilon_i,\quad\epsilon_i\sim N_2(0,\Sigma_j),\label{q15:c0}
\end{equation}
a SUR system. Prior times likelihood:
\begin{equation}
\pi(\beta\mid\cdot)\propto\exp\left\{-\tfrac12(\beta-b_0)'B_0^{-1}(\beta-b_0)\right\}\prod_{j=0}^{1}\prod_{i\in N_j}\exp\left\{-\tfrac12(y_i^*-X_{ji}\beta)'\Sigma_j^{-1}(y_i^*-X_{ji}\beta)\right\}.\label{q15:c1}
\end{equation}
Multiplying the exponent of \eqref{q15:c1} by $-2$ and collecting terms in $\beta$:
\begin{equation}
\beta'\Big(B_0^{-1}+\sum_{j}\sum_{i\in N_j}X_{ji}'\Sigma_j^{-1}X_{ji}\Big)\beta-2\beta'\Big(B_0^{-1}b_0+\sum_{j}\sum_{i\in N_j}X_{ji}'\Sigma_j^{-1}y_i^*\Big)+\text{const}.\label{q15:c2}
\end{equation}
Writing the first bracket as $B^{-1}$ and the second as $B^{-1}\hat\beta$, \eqref{q15:c2} becomes
\begin{equation}
(\beta-\hat\beta)'B^{-1}(\beta-\hat\beta)+\text{const},\label{q15:c3}
\end{equation}
so $\beta\mid\cdot\sim N_p(\hat\beta,B)$ with $B$ and $\hat\beta$ exactly as in step 11.2.}

\passo{(d) $s_i^*$}{Given everything else, $s_i^*$ enters only unit $i$'s joint density and the indicator
that ties it to the observed $s_i$:
\begin{equation}
\pi(s_i^*\mid\cdot)\propto N_2\big(y_i^*\mid X_{ji}\beta,\Sigma_j\big)\big[1(s_i=0)1(s_i^*\leq0)+1(s_i=1)1(s_i^*>0)\big].\label{q15:d0}
\end{equation}
Factor the bivariate density as $f(y_{ji})f(s_i^*\mid y_{ji})$; $f(y_{ji})$ does not involve $s_i^*$, and by
\eqref{q14:mu}--\eqref{q14:phi} of Question 14 $f(s_i^*\mid y_{ji})=N(s_i^*\mid\mu_{ji},\phi_j^2)$. Substituting
into \eqref{q15:d0}:
\begin{equation}
\pi(s_i^*\mid\cdot)\propto N\big(s_i^*\mid\mu_{ji},\phi_j^2\big)\big[1(s_i=0)1(s_i^*\leq0)+1(s_i=1)1(s_i^*>0)\big],\label{q15:d1}
\end{equation}
i.e.\ $\mathrm{TN}_{(-\infty,0)}(\mu_{0i},\phi_0^2)$ when $s_i=0$ and $\mathrm{TN}_{(0,\infty)}(\mu_{1i},\phi_1^2)$ when
$s_i=1$, as in step 11.3 (whether $0$ is included is a set of probability zero).}

\intuicao{Every block is conjugate once the latent $s^*$ is filled in: given $s^*$ the model is a normal
SUR regression with a control-function term, so the Gibbs sampler only cycles through inverse-gamma,
normal and truncated-normal draws.}
\refbib{Greenberg \S11.1 (pp.~170--172), Algorithm 11.1; \S9.1 (pp.~134--139), SUR; \S8.2.2, truncated-normal latent data.}
\end{sol}

\end{document}
